% =====================================================================
%  Respuesta en frecuencia y compensación en frecuencia
%  Apunte de estudio — Sistemas de Control
%  Compilar con:  pdflatex apunte.tex  (dos pasadas)  o  latexmk -pdf apunte.tex
%  Las figuras leen los datos de la carpeta datos/ (generar_datos.py)
% =====================================================================
\documentclass[11pt,a4paper]{article}
\usepackage[utf8]{inputenc}
\usepackage[T1]{fontenc}
\usepackage{lmodern}
\usepackage[spanish,es-noshorthands,es-tabla]{babel}
\usepackage{amsmath,amssymb}
\usepackage{icomma}
\usepackage[a4paper,left=2.2cm,right=2.2cm,top=2.3cm,bottom=2.3cm,headheight=14pt]{geometry}
\usepackage{xcolor}
\usepackage{booktabs,array,tabularx}
\usepackage{enumitem}
\usepackage{needspace}
\usepackage{titlesec}
\usepackage{fancyhdr}
\usepackage[most]{tcolorbox}
\usepackage{pgfplots}
\pgfplotsset{compat=1.18}
\usepgfplotslibrary{groupplots}
\usetikzlibrary{arrows.meta,positioning,calc,decorations.markings}
\usepackage[font=small,labelfont=bf]{caption}
\usepackage[colorlinks=true,linkcolor=azul,urlcolor=azul]{hyperref}

% ---------------------------------------------------------------- colores
\definecolor{azul}{RGB}{31,78,121}
\definecolor{verde}{RGB}{0,115,70}
\definecolor{rojo}{RGB}{185,30,40}
\definecolor{naranja}{RGB}{215,110,0}
\definecolor{violeta}{RGB}{110,50,150}
\definecolor{gris}{RGB}{95,95,95}

% ---------------------------------------------------------------- títulos
\titleformat{\section}{\Large\bfseries\color{azul}}{\thesection.}{0.5em}{}
\titleformat{\subsection}{\large\bfseries\color{azul}}{\thesubsection}{0.5em}{}
\titleformat{\subsubsection}{\normalsize\bfseries\color{azul}}{}{0em}{}
\titlespacing*{\section}{0pt}{14pt plus 3pt}{6pt}
\titlespacing*{\subsection}{0pt}{10pt plus 2pt}{4pt}
\titlespacing*{\subsubsection}{0pt}{8pt plus 2pt}{3pt}
\setlength{\parindent}{0pt}
\setlength{\parskip}{4pt plus 1pt}
\setlist{itemsep=1pt,topsep=2pt,parsep=1pt}

\pagestyle{fancy}
\fancyhf{}
\lhead{\small\color{gris}Sistemas de Control}
\rhead{\small\color{gris}Respuesta en frecuencia y compensación}
\cfoot{\small\thepage}
\renewcommand{\headrulewidth}{0.4pt}

% ---------------------------------------------------------------- cajas
\tcbset{base/.style={enhanced,breakable,boxrule=0.6pt,arc=2pt,left=6pt,right=6pt,top=4pt,bottom=4pt,
  fonttitle=\bfseries\small,before skip=6pt,after skip=6pt}}
\newtcolorbox{idea}[1][Idea clave]{base,colback=azul!5,colframe=azul,title={#1}}
\newtcolorbox{receta}[1][Procedimiento]{base,colback=verde!5,colframe=verde,title={#1}}
\newtcolorbox{cuidado}[1][Ojo]{base,colback=rojo!4,colframe=rojo,title={#1}}
\newtcolorbox{formula}{base,colback=azul!4,colframe=azul!60,boxrule=0.5pt}
\newcounter{ejemplo}
\newtcolorbox{ejemplobox}[1]{base,colback=gris!4,colframe=gris,title={Ejemplo \theejemplo\ --- #1}}
\newenvironment{ejemplo}[1]{\refstepcounter{ejemplo}\begin{ejemplobox}{#1}}{\end{ejemplobox}}
\setcounter{tocdepth}{1}
\setcounter{secnumdepth}{2}

% ejercicios y soluciones
\newcounter{ejer}
\newcommand{\ejercicio}[1]{\par\medskip\needspace{5\baselineskip}\refstepcounter{ejer}\noindent{\color{azul}\bfseries Ejercicio \theejer\ \ }{\bfseries #1}\par\nopagebreak\smallskip}
\newcommand{\solucion}[2]{\par\medskip\needspace{5\baselineskip}\noindent{\color{verde}\bfseries Solución \ref{#1}\ \ }{\bfseries #2}\par\nopagebreak\smallskip}
\newcommand{\rta}[1]{\fbox{$\displaystyle #1$}}

% partes
\newcommand{\parte}[2]{\clearpage
  \begin{tcolorbox}[enhanced,colback=azul,colframe=azul,arc=3pt,left=10pt,top=8pt,bottom=8pt]
  \color{white}\Large\bfseries Parte #1\hfill\normalsize #2\end{tcolorbox}\vspace{2pt}}

% ---------------------------------------------------------------- macros
\newcommand{\pct}{\ensuremath{\,\mathchar"7025}}
\newcommand{\dB}{\,\mathrm{dB}}
\newcommand{\dBdec}{\,\mathrm{dB/dec}}
\newcommand{\rads}{\,\mathrm{rad/s}}
\newcommand{\jw}{j\omega}
\newcommand{\MF}{\mathrm{MF}}
\newcommand{\MG}{\mathrm{MG}}
\newcommand{\wc}{\omega_c}
\newcommand{\wf}{\omega_{-180}}

% ---------------------------------------------------------------- pgfplots
\pgfplotsset{
  every axis/.append style={
    tick label style={font=\scriptsize}, label style={font=\footnotesize},
    legend style={font=\scriptsize,fill=white,fill opacity=0.9,text opacity=1,draw=gray!50},
    grid=both, major grid style={gray!28}, minor grid style={gray!10}},
  bode/.style={xmode=log, log basis x=10, minor x tick num=8},
  exacta/.style={azul, thick},
  asint/.style={rojo, thick, dashed},
  antes/.style={gris, thick, dashed},
  despues/.style={azul, thick},
  comp/.style={naranja, thick, dash dot},
}

\pgfkeys{/pgf/number format/.cd,use comma,1000 sep={}}

\begin{document}

% =====================================================================
%  PORTADA
% =====================================================================
\begin{center}
{\color{azul}\Huge\bfseries Respuesta en frecuencia\\[2pt] y compensación en frecuencia}\\[8pt]
{\large Apunte de estudio con ejemplos resueltos y ejercicios}\\[4pt]
{\normalsize Sistemas de Control}
\end{center}
\vspace{2pt}

\begin{idea}[Cómo usar este apunte]
La \textbf{Parte I} desarrolla las curvas de respuesta en frecuencia: Bode, diagrama polar, Nyquist,
Nichols, márgenes de estabilidad y la correlación con la respuesta temporal. La \textbf{Parte II} usa esas
herramientas para diseñar compensadores de \textbf{adelanto}, \textbf{atraso} y \textbf{atraso-adelanto}.
Cada tema tiene un ejemplo resuelto con números verificados por simulación.
Los ejercicios están al final de cada parte y las \textbf{soluciones completas} en la sección~\ref{sec:soluciones}.
Conviene intentarlos antes de mirar la solución. El formulario del Apéndice~\ref{ap:formulario} resume todo en una página
y el Apéndice~\ref{ap:matlab} muestra cómo verificar cada resultado en MATLAB.
\end{idea}

{\small\tableofcontents}

% =====================================================================
\parte{I}{Curvas de respuesta en frecuencia}
% =====================================================================

\section{Qué es la respuesta en frecuencia}

\subsection{Régimen senoidal permanente}
Sea un sistema lineal, invariante y \textbf{estable} con transferencia $G(s)$. Si la entrada es
$u(t)=A\sin(\omega t)$, una vez extinguido el transitorio la salida es otra senoide de la
\emph{misma frecuencia}:
\begin{formula}
\[
u(t)=A\sin(\omega t)\quad\Longrightarrow\quad y_{ss}(t)=A\,|G(\jw)|\,\sin\!\big(\omega t+\angle G(\jw)\big).
\]
\end{formula}
La justificación es corta. $U(s)=A\omega/(s^2+\omega^2)$, y al expandir $Y(s)=G(s)U(s)$ en fracciones simples
los términos de los polos de $G$ decaen (porque $G$ es estable). Solo sobreviven los términos de los polos $\pm\jw$
de la entrada, cuyos residuos dependen de $G(\pm\jw)$. Al sumarlos queda la expresión de arriba.

El número complejo $G(\jw)$ se llama \textbf{respuesta en frecuencia}. Para cada $\omega$, el módulo $|G(\jw)|$
es la \textbf{ganancia} (cuánto se amplifica la senoide) y el argumento $\angle G(\jw)$ es el
\textbf{desfasaje} (cuánto se adelanta o atrasa). Para obtenerla alcanza con reemplazar $s\to\jw$.

\subsection{Por qué conviene trabajar en frecuencia}
\begin{itemize}
\item \textbf{Se puede medir.} Con un generador de senoides y un osciloscopio se obtiene $G(\jw)$
      aunque no se conozca el modelo. Es la base de la identificación experimental.
\item \textbf{Se diseña sin calcular raíces.} La estabilidad y el amortiguamiento del lazo cerrado se leen
      sobre la respuesta del \emph{lazo abierto} $L(\jw)=G_c(\jw)G(\jw)H(\jw)$.
\item \textbf{Los productos se vuelven sumas.} En escala logarítmica, conectar bloques en serie es sumar
      curvas. Por eso agregar un compensador se entiende y se diseña gráficamente.
\item \textbf{Retardos, ruido y robustez.} Un retardo puro $e^{-Ts}$ no es racional y complica el lugar de raíces,
      pero en frecuencia es solo un término de fase. El ruido de medición y la dinámica no modelada viven en alta
      frecuencia, y ahí se los ve directamente.
\end{itemize}

\subsection{Tres maneras de dibujar $G(\jw)$}
\begin{center}
\small
\begin{tabular}{@{}lll@{}}
\toprule
Gráfico & Ejes & Uso principal\\
\midrule
Bode (logarítmico) & $20\log|G|$ y $\angle G$ vs.\ $\log\omega$ (dos curvas) & Trazado a mano, diseño de compensadores\\
Polar / Nyquist & $\mathrm{Re}\,G$ vs.\ $\mathrm{Im}\,G$, con $\omega$ como parámetro & Criterio de estabilidad de Nyquist\\
Nichols (mag.--fase) & $20\log|G|$ vs.\ $\angle G$ & Leer el lazo cerrado sobre el lazo abierto\\
\bottomrule
\end{tabular}
\end{center}
Las tres muestran la misma información. Cambia qué propiedad se lee más fácil en cada una.

% ---------------------------------------------------------------------
\section{Diagrama de Bode}

\subsection{Decibeles, décadas y octavas}
La magnitud se expresa en decibeles: $|G|_{\mathrm{dB}}=20\log_{10}|G|$. La frecuencia va en escala logarítmica.
Una \textbf{década} es un factor 10 en $\omega$ y una \textbf{octava} es un factor 2.
\begin{center}
\small
\begin{tabular}{@{}l*{9}{c}@{}}
\toprule
$|G|$ & 0,01 & 0,1 & 0,5 & $1/\sqrt2$ & 1 & $\sqrt2$ & 2 & 10 & 100\\
dB & $-40$ & $-20$ & $-6{,}02$ & $-3{,}01$ & 0 & $+3{,}01$ & $+6{,}02$ & $+20$ & $+40$\\
\bottomrule
\end{tabular}
\end{center}
Una pendiente de $-20\dBdec$ equivale a $-6\,\mathrm{dB/octava}$. Las magnitudes se multiplican, así que
sus dB \textbf{se suman}. Las fases también se suman.

\subsection{Forma de Bode (de constantes de tiempo)}
Antes de graficar, se reescribe la transferencia de modo que cada factor valga 1 en baja frecuencia:
\begin{formula}
\[
G(s)=\frac{K_B\displaystyle\prod_i\Big(1+\frac{s}{z_i}\Big)\prod_l\Big(\frac{s^2}{\omega_{nl}'^2}+\frac{2\zeta_l'}{\omega_{nl}'}s+1\Big)}
{s^N\displaystyle\prod_k\Big(1+\frac{s}{p_k}\Big)\prod_m\Big(\frac{s^2}{\omega_{nm}^2}+\frac{2\zeta_m}{\omega_{nm}}s+1\Big)}
\qquad\text{con}\quad K_B=\lim_{s\to0}s^N G(s).
\]
\end{formula}
$N$ es el número de integradores (\textbf{tipo} del sistema) y $K_B$ es la \textbf{ganancia de Bode}.
Para pasar de la forma polo-cero a esta forma se saca factor común el término independiente de cada factor.
Por ejemplo, $\frac{10(s+2)}{s(s+5)}=\frac{10\cdot2}{5}\,\frac{(1+s/2)}{s(1+s/5)}$, de modo que $K_B=4$.

\subsection{Factores elementales}

\subsubsection{a) Ganancia constante $K_B$}
Magnitud plana en $20\log|K_B|$. La fase es $0^\circ$ si $K_B>0$ y $-180^\circ$ si $K_B<0$.

\subsubsection{b) Polos y ceros en el origen: $(\jw)^{\mp N}$}
La magnitud es $\mp 20N\log\omega$, una recta de pendiente $\mp20N\dBdec$ que pasa por $0\dB$ en $\omega=1$.
La fase es constante, $\mp90^\circ N$.

\subsubsection{c) Factor de primer orden: $(1+\jw/\omega_b)^{\mp1}$}
Para el polo $1/(1+\jw/\omega_b)$:
\[
|G|_{\mathrm{dB}}=-10\log\!\big(1+(\omega/\omega_b)^2\big),\qquad \angle G=-\arctan(\omega/\omega_b).
\]
\begin{itemize}
\item \textbf{Asíntotas.} Para $\omega\ll\omega_b$ la magnitud es $0\dB$; para $\omega\gg\omega_b$ es una recta de
      $-20\dBdec$. Se cortan en la \textbf{frecuencia de quiebre} $\omega_b$.
\item \textbf{Error.} El error de la asíntota es máximo en el quiebre ($-3\dB$) y vale $-1\dB$ una octava antes o después.
\item \textbf{Fase.} Vale $-45^\circ$ en $\omega_b$. Se aproxima por una recta que va de $0^\circ$ en
      $\omega_b/10$ a $-90^\circ$ en $10\,\omega_b$, con error máximo de $5{,}7^\circ$.
\item \textbf{Cero.} El cero $(1+\jw/\omega_b)$ es la imagen especular: $+20\dBdec$ y fase de $0^\circ$ a $+90^\circ$.
\end{itemize}

\begin{figure}[ht]
\centering
\begin{tikzpicture}
\begin{groupplot}[group style={group size=2 by 1,horizontal sep=1.6cm},
  bode,width=0.48\linewidth,height=4.6cm,xmin=0.01,xmax=100,xlabel={$\omega/\omega_b$}]
\nextgroupplot[ylabel={Magnitud [dB]},ymin=-42,ymax=5,ytick={-40,-30,-20,-10,-3,0},
  legend pos=south west]
\addplot[exacta] table[x=u,y=mag]{datos/primer_orden.dat};
\addplot[asint] table[x=u,y=asin]{datos/primer_orden.dat};
\legend{Exacta,Asintótica}
\node[font=\scriptsize,anchor=south west] at (axis cs:1.1,-3) {$-3\dB$};
\node[font=\scriptsize,rotate=-27] at (axis cs:4,-24) {$-20\dBdec$};
\nextgroupplot[ylabel={Fase [$^\circ$]},ymin=-95,ymax=5,ytick={-90,-45,0}]
\addplot[exacta] table[x=u,y=fase]{datos/primer_orden.dat};
\addplot[asint] table[x=u,y=fasin]{datos/primer_orden.dat};
\end{groupplot}
\end{tikzpicture}
\caption{Polo de primer orden $1/(1+\jw/\omega_b)$: curva exacta y aproximación por rectas.}
\label{fig:primer}
\end{figure}

\subsubsection{d) Factor de segundo orden: $\big[1-(\omega/\omega_n)^2+j2\zeta\,\omega/\omega_n\big]^{-1}$}
Las asíntotas son $0\dB$ hasta $\omega_n$ y $-40\dBdec$ después. La curva real depende mucho de $\zeta$:
\[
|G(j\omega_n)|=\frac{1}{2\zeta}\ \Rightarrow\ \text{error en el quiebre}=-20\log(2\zeta),\qquad \angle G(j\omega_n)=-90^\circ.
\]
Si $\zeta<1/\sqrt2\approx0{,}707$ aparece un pico de \textbf{resonancia}:
\begin{formula}
\[
M_r=\frac{1}{2\zeta\sqrt{1-\zeta^2}},\qquad \omega_r=\omega_n\sqrt{1-2\zeta^2}.
\]
\end{formula}
La fase va de $0^\circ$ a $-180^\circ$, y la transición es más abrupta cuanto menor es $\zeta$ (figura~\ref{fig:segundo}).
La aproximación asintótica es razonable para $0{,}4\lesssim\zeta\le1$. Fuera de ese rango hay que corregir
el pico a mano con $-20\log(2\zeta)$.

\begin{figure}[ht]
\centering
\begin{tikzpicture}
\begin{groupplot}[group style={group size=2 by 1,horizontal sep=1.6cm},
  bode,width=0.48\linewidth,height=5cm,xmin=0.1,xmax=10,xlabel={$\omega/\omega_n$}]
\nextgroupplot[ylabel={Magnitud [dB]},ymin=-40,ymax=22,legend pos=south west,legend columns=2]
\addplot[violeta,thick] table[x=u,y=m005]{datos/segundo_orden.dat};
\addplot[rojo,thick] table[x=u,y=m01]{datos/segundo_orden.dat};
\addplot[naranja,thick] table[x=u,y=m02]{datos/segundo_orden.dat};
\addplot[verde,thick] table[x=u,y=m03]{datos/segundo_orden.dat};
\addplot[azul,thick] table[x=u,y=m05]{datos/segundo_orden.dat};
\addplot[cyan!70!black,thick] table[x=u,y=m0707]{datos/segundo_orden.dat};
\addplot[black,thick] table[x=u,y=m1]{datos/segundo_orden.dat};
\addplot[black,dashed] table[x=u,y=asin]{datos/segundo_orden.dat};
\legend{$\zeta{=}0{,}05$,$0{,}1$,$0{,}2$,$0{,}3$,$0{,}5$,$0{,}707$,$1$}
\nextgroupplot[ylabel={Fase [$^\circ$]},ymin=-185,ymax=5,ytick={-180,-135,-90,-45,0}]
\addplot[violeta,thick] table[x=u,y=f005]{datos/segundo_orden.dat};
\addplot[rojo,thick] table[x=u,y=f01]{datos/segundo_orden.dat};
\addplot[naranja,thick] table[x=u,y=f02]{datos/segundo_orden.dat};
\addplot[verde,thick] table[x=u,y=f03]{datos/segundo_orden.dat};
\addplot[azul,thick] table[x=u,y=f05]{datos/segundo_orden.dat};
\addplot[cyan!70!black,thick] table[x=u,y=f0707]{datos/segundo_orden.dat};
\addplot[black,thick] table[x=u,y=f1]{datos/segundo_orden.dat};
\end{groupplot}
\end{tikzpicture}
\caption{Factor de segundo orden para distintos $\zeta$. La línea negra discontinua es la asíntota.}
\label{fig:segundo}
\end{figure}

\subsubsection{e) Retardo puro: $e^{-\jw T}$}
$|e^{-\jw T}|=1$, es decir $0\dB$: \textbf{no modifica la magnitud}. Su fase es $-\omega T$ (en radianes), o sea
$-57{,}3\,\omega T$ grados. Crece linealmente con $\omega$, y en el eje logarítmico se ve como una caída cada
vez más pronunciada. Es el peor enemigo de la estabilidad: resta fase sin dar ninguna señal en la magnitud.

\begin{center}
\small
\begin{tabular}{@{}lccc@{}}
\toprule
Factor & Quiebre & Cambio de pendiente & Fase aportada (total)\\
\midrule
$K_B$ & --- & --- & $0^\circ$ ($-180^\circ$ si $K_B<0$)\\
$1/s^N$ & --- & $-20N\dBdec$ en todo $\omega$ & $-90^\circ N$\\
Polo real $1/(1+s/p)$ & $p$ & $-20\dBdec$ & $0\to-90^\circ$\\
Cero real $(1+s/z)$ & $z$ & $+20\dBdec$ & $0\to+90^\circ$\\
Par de polos complejos & $\omega_n$ & $-40\dBdec$ & $0\to-180^\circ$\\
Cero de fase no mínima $(1-s/z)$ & $z$ & $+20\dBdec$ & $0\to-90^\circ$\\
Retardo $e^{-Ts}$ & --- & ninguno & $-\omega T$ (sin límite)\\
\bottomrule
\end{tabular}
\end{center}

\subsection{Receta para trazar el Bode asintótico}
\begin{receta}[Trazado a mano en 6 pasos]
\begin{enumerate}[label=\textbf{\arabic*.}]
\item Escribir $G$ en forma de Bode y anotar $K_B$, $N$ y las frecuencias de quiebre ordenadas de menor a mayor.
\item \textbf{Tramo inicial.} La recta $20\log K_B-20N\log\omega$, que vale $20\log K_B$ en $\omega=1$
      y tiene pendiente $-20N\dBdec$.
\item Recorrer los quiebres de izquierda a derecha. En cada uno la pendiente cambia: $\pm20\dBdec$ en cada cero o polo
      real y $\pm40\dBdec$ en cada par complejo. Un quiebre doble cuenta dos veces.
\item Verificar: la pendiente final es $-20(n-m)\dBdec$ y la fase final es $-90^\circ(n-m)$, si el sistema es de fase mínima.
\item Corregir la curva: $\pm3\dB$ en los quiebres simples y $-20\log(2\zeta)$ en los pares complejos.
\item \textbf{Fase.} Sumar las contribuciones de cada factor en algunos puntos clave (los quiebres, $\wc$)
      con $\arctan$, o con las aproximaciones por rectas.
\end{enumerate}
\end{receta}

\begin{ejemplo}{Trazado completo}
\label{ej:bode}
Trazar el Bode de $\displaystyle G(s)=\frac{16000\,(s+1)}{s\,(s+5)\,(s^2+16s+1600)}$.

\textbf{Forma de Bode.} $s^2+16s+1600=1600\big[(s/40)^2+0{,}01s+1\big]$, de modo que $\omega_n=40$ y
$2\zeta/\omega_n=0{,}01$, con lo cual $\zeta=0{,}2$. Sacando factor común:
\[
G(s)=\frac{16000\cdot1}{5\cdot1600}\cdot\frac{(1+s)}{s\,(1+s/5)\big[(s/40)^2+0{,}4\,(s/40)+1\big]}
=\frac{2\,(1+s)}{s\,(1+s/5)\big[(s/40)^2+0{,}4\,(s/40)+1\big]}.
\]
$K_B=2$, $N=1$. Quiebres: cero en 1, polo en 5 y par complejo en 40 con $\zeta=0{,}2$.

\textbf{Asíntotas.}
\begin{itemize}
\item $\omega<1$: recta $20\log(2/\omega)$, de $-20\dBdec$. En $\omega=1$ vale $6{,}02\dB$.
\item $1<\omega<5$: el cero lleva la pendiente a $0\dBdec$ y la curva queda plana en $6{,}02\dB$.
\item $5<\omega<40$: el polo la lleva a $-20\dBdec$. En $\omega=40$ vale $6{,}02-20\log 8=-12{,}0\dB$.
      Cruza $0\dB$ en $6{,}02-20\log(\omega/5)=0$, es decir en $\omega=10\rads$.
\item $\omega>40$: el par complejo agrega $-40$ y la pendiente final es $-60\dBdec$. Verificación: $n-m=4-1=3$.
\end{itemize}
\textbf{Correcciones.} $+3\dB$ en $\omega=1$, $-3\dB$ en $\omega=5$ y $-20\log(0{,}4)=+8\dB$ en $\omega=40$. El pico real
queda cerca de $-4\dB$. La simulación da $-3{,}4\dB$ en $\omega=36{,}3\rads$, desplazado por la caída de $-20\dBdec$ que lo precede.

\textbf{Fase.} Arranca en $-90^\circ$ (integrador), sube por el cero, baja por el polo, cae bruscamente cerca de 40
y termina en $-90^\circ\cdot3=-270^\circ$. En $\omega=10$:
$-90+\arctan10-\arctan2-\arctan\frac{0{,}4\cdot0{,}25}{1-0{,}25^2}=-90+84{,}3-63{,}4-6{,}1=-75{,}2^\circ$.

\textbf{Dato curioso.} El cruce exacto es $\wc=9{,}33\rads$ con fase $\approx-73{,}6^\circ$, así que el margen de fase es enorme ($106^\circ$).
Sin embargo, el pico resonante queda a solo $4{,}5\dB$ de $0\dB$ justo donde la fase pasa por $-180^\circ$
($\omega=40{,}8\rads$). El margen de ganancia es apenas $4{,}5\dB$: \emph{un margen bueno no garantiza el otro}.
\end{ejemplo}

\begin{figure}[ht]
\centering
\begin{tikzpicture}
\begin{groupplot}[group style={group size=1 by 2,vertical sep=0.35cm,xticklabels at=edge bottom},
  bode,width=0.92\linewidth,height=4.6cm,xmin=0.01,xmax=1000]
\nextgroupplot[ylabel={$|G|$ [dB]},ymin=-110,ymax=50,ytick={-100,-80,-60,-40,-20,0,20,40},legend pos=south west]
\addplot[exacta] table[x=w,y=mag]{datos/bode_ejemplo.dat};
\addplot[asint] table[x=w,y=asin]{datos/bode_ejemplo.dat};
\legend{Exacta,Asintótica}
\draw[gris,densely dotted] (axis cs:1,-110)--(axis cs:1,50) (axis cs:5,-110)--(axis cs:5,50) (axis cs:40,-110)--(axis cs:40,50);
\node[font=\scriptsize,rojo] at (axis cs:0.06,36) {$-20$};
\node[font=\scriptsize,rojo] at (axis cs:2.2,14) {$0$};
\node[font=\scriptsize,rojo] at (axis cs:12,4) {$-20$};
\node[font=\scriptsize,rojo] at (axis cs:250,-50) {$-60\dBdec$};
\nextgroupplot[ylabel={$\angle G$ [$^\circ$]},xlabel={$\omega$ [rad/s]},ymin=-280,ymax=0,ytick={-270,-180,-90,0}]
\addplot[exacta] table[x=w,y=fase]{datos/bode_ejemplo.dat};
\draw[gris,densely dotted] (axis cs:1,-280)--(axis cs:1,0) (axis cs:5,-280)--(axis cs:5,0) (axis cs:40,-280)--(axis cs:40,0);
\draw[black,dashed] (axis cs:0.01,-180)--(axis cs:1000,-180);
\end{groupplot}
\end{tikzpicture}
\caption{Bode del ejemplo~\ref{ej:bode}. Las líneas punteadas marcan los quiebres (1, 5 y 40~rad/s).}
\label{fig:bodeej}
\end{figure}

\subsection{Tipo del sistema y constantes de error leídas del Bode}
El tramo de baja frecuencia del Bode de lazo abierto dice cuánto error tendrá el lazo cerrado en régimen permanente:
\begin{center}
\small
\begin{tabular}{@{}lllll@{}}
\toprule
Tipo & Pendiente & Constante & Cómo leerla en el Bode & Error (realim.\ unitaria)\\
\midrule
0 & $0\dBdec$ & $K_p=\lim_{s\to0}L(s)$ & altura del tramo plano: $20\log K_p$ & escalón: $1/(1+K_p)$\\
1 & $-20\dBdec$ & $K_v=\lim_{s\to0}sL(s)$ & asíntota inicial en $0\dB$: $\omega=K_v$ & rampa: $1/K_v$\\
2 & $-40\dBdec$ & $K_a=\lim_{s\to0}s^2L(s)$ & asíntota inicial en $0\dB$: $\omega=\sqrt{K_a}$ & parábola: $1/K_a$\\
\bottomrule
\end{tabular}
\end{center}
\begin{cuidado}
La constante se lee sobre la \textbf{prolongación de la asíntota inicial}, no sobre la curva real. Si hay un
cero o un polo antes de $\omega=K_v$, la curva real ya se separó de la asíntota.
En el ejemplo~\ref{ej:bode}, $K_v=K_B=2$: la asíntota $2/\omega$ corta $0\dB$ en $\omega=2$, aunque la curva real cruce en 9,33.
\end{cuidado}

\subsection{Fase mínima, fase no mínima e identificación}
Un sistema es de \textbf{fase mínima} si no tiene polos ni ceros en el semiplano derecho y no tiene retardos.
Para estos sistemas, Bode demostró que \emph{la magnitud determina la fase}. En un tramo de pendiente constante
$-20k\dBdec$ que se extiende alrededor de una década, la fase se acerca a $-90^\circ k$.

Un cero $(1-s/z)$ tiene \emph{la misma magnitud} que $(1+s/z)$ pero resta fase en lugar de sumarla.
En el Bode, un sistema de fase no mínima se reconoce porque tiene más atraso de fase del que su pendiente justifica.
Estos sistemas son más difíciles de controlar: la respuesta al escalón arranca para el lado contrario (subestimación)
y el cruce no puede llevarse muy alto.

\begin{receta}[Identificar $G(s)$ a partir de un Bode medido]
\begin{enumerate}[label=\textbf{\arabic*.}]
\item Trazar rectas de pendientes múltiplos de $\pm20\dBdec$ que se ajusten a la curva medida.
\item La pendiente inicial da el tipo $N$. La altura del tramo inicial da $K_B$ (tabla anterior).
\item Cada cambio de pendiente es un quiebre: $-20$ indica un polo, $+20$ un cero y $-40$ un par de polos
      (o un polo doble).
\item En un cambio de $-40$, la distancia de la curva real a la asíntota en el quiebre indica el tipo de par:
      unos $-6\dB$ corresponden a un polo real doble ($\zeta=1$) y un pico corresponde a $\zeta=10^{-\Delta/20}/2$.
\item Comparar la fase medida con la que predice el modelo. Si sobra atraso, hay un retardo o un cero de fase no mínima.
\end{enumerate}
\end{receta}

% ---------------------------------------------------------------------
\section{Diagrama polar y criterio de Nyquist}

\subsection{Diagrama polar}
Es la curva que recorre $G(\jw)$ en el plano complejo cuando $\omega$ va de $0$ a $\infty$. Su forma
general se deduce del comportamiento en los dos extremos:
\begin{itemize}
\item $\omega\to0$: lo determina el tipo. Un sistema tipo 0 arranca en $K_p$ sobre el eje real positivo.
      Uno tipo 1 arranca en el infinito en dirección $-90^\circ$, con una asíntota vertical en
      $\mathrm{Re}=-K_v\sum T_i$ ($T_i$: constantes de tiempo de los polos menos las de los ceros).
      Uno tipo 2 arranca en el infinito en dirección $-180^\circ$.
\item $\omega\to\infty$: si $n>m$, llega al origen tangente a la dirección $-90^\circ(n-m)$.
\end{itemize}

\begin{figure}[ht]
\centering
\begin{tikzpicture}
\begin{groupplot}[group style={group size=3 by 1,horizontal sep=1.1cm},
  width=0.36\linewidth,height=0.36\linewidth,axis equal image,grid=major,
  xlabel={Re},ylabel={Im},title style={font=\footnotesize}]
\nextgroupplot[title={Tipo 0: $\frac{1}{(s+1)(0{,}5s+1)(0{,}2s+1)}$},xmin=-0.6,xmax=1.2,ymin=-1,ymax=0.6]
\addplot[exacta] table[x=re,y=im]{datos/polar_t0.dat};
\addplot[only marks,mark=*,mark size=1.3pt] coordinates {(1,0)};
\node[font=\scriptsize,anchor=south] at (axis cs:1,0.02) {$\omega{=}0$};
\nextgroupplot[title={Tipo 1: $\frac{1}{s(s+1)(0{,}5s+1)}$},xmin=-2.4,xmax=0.6,ymin=-3,ymax=0.6]
\addplot[exacta] table[x=re,y=im]{datos/polar_t1.dat};
\draw[rojo,dashed] (axis cs:-1.5,-3)--(axis cs:-1.5,0.6);
\node[font=\scriptsize,rojo,anchor=west] at (axis cs:-1.45,-2.6) {$-1{,}5$};
\nextgroupplot[title={Tipo 2: $\frac{1}{s^2(s+1)}$},xmin=-3,xmax=0.5,ymin=-0.5,ymax=3]
\addplot[exacta] table[x=re,y=im]{datos/polar_t2.dat};
\end{groupplot}
\end{tikzpicture}
\caption{Formas típicas del diagrama polar ($\omega$ de 0 a $\infty$). Todas terminan en el origen con dirección $-270^\circ$ porque $n-m=3$.}
\label{fig:polar}
\end{figure}

Los cruces con los ejes se calculan igualando a cero la parte real o la imaginaria.
\textbf{El cruce con el eje real negativo es el dato más importante}, porque decide la estabilidad.

\subsection{El criterio de Nyquist}
En un lazo con realimentación negativa, los polos de lazo cerrado son las raíces de $1+L(s)=0$.
El criterio cuenta cuántas de esas raíces están en el semiplano derecho (SPD) sin calcularlas.

\textbf{Principio del argumento.} Si un contorno cerrado del plano $s$ encierra en sentido horario $Z$ ceros y $P$
polos de una función $F(s)$, la imagen $F(s)$ rodea al origen $N=Z-P$ veces en sentido horario.

\textbf{Aplicación.} Se toma $F=1+L$ y como contorno el \textbf{contorno de Nyquist}, que encierra todo el SPD:
el eje imaginario completo más un semicírculo de radio infinito. Los polos de $1+L$ son los de $L$, y los ceros de $1+L$
son los polos de lazo cerrado. Rodear el origen con $1+L$ equivale a rodear el punto $-1$ con $L$.

\begin{formula}
\[
\boxed{Z=N+P}\qquad
\begin{array}{l}
Z:\ \text{polos de lazo cerrado en el SPD (se busca }Z=0)\\
P:\ \text{polos de lazo abierto }L(s)\text{ en el SPD (dato)}\\
N:\ \text{rodeos de }L(\jw)\text{ al punto }-1\text{ (horario }+,\ \text{antihorario }-)
\end{array}
\]
\end{formula}

\begin{receta}[Aplicar Nyquist]
\begin{enumerate}[label=\textbf{\arabic*.}]
\item Contar $P$, los polos de $L(s)$ con parte real positiva.
\item Trazar $L(\jw)$ para $\omega:0\to\infty$ y su espejo respecto del eje real ($\omega<0$).
\item Si $L$ tiene polos sobre el eje imaginario (integradores), el contorno los esquiva por la derecha con un semicírculo
      de radio $\varepsilon\to0$. Cada integrador se transforma en una semicircunferencia de radio infinito
      recorrida en \emph{sentido horario}: $q$ integradores dan $q\times180^\circ$.
\item Contar $N$: trazar una semirrecta desde $-1$ hacia cualquier lado y sumar los cruces de la curva
      ($+1$ horario, $-1$ antihorario).
\item $Z=N+P$. El lazo cerrado es estable si y solo si $Z=0$.
\end{enumerate}
\end{receta}

\begin{idea}[Caso más común: lazo abierto estable ($P=0$)]
Si $L(s)$ no tiene polos en el SPD, el lazo cerrado es estable si y solo si \textbf{la curva de Nyquist no rodea al punto $-1$}.
Para sistemas de fase mínima esto equivale a pedir que, en la frecuencia donde la fase vale $-180^\circ$, el módulo sea menor que 1.
De ahí salen los márgenes de la sección~\ref{sec:margenes}.
\end{idea}

\begin{ejemplo}{Nyquist y ganancia crítica}
\label{ej:nyq}
Sea $L(s)=\dfrac{K}{s(s+1)(s+2)}$, con $K>0$ y $P=0$. ¿Para qué $K$ el lazo cerrado es estable?

\textbf{Cruce con el eje real.} El denominador de $L(\jw)$ es $\jw(\jw+1)(\jw+2)=-3\omega^2+j\omega(2-\omega^2)$.
Es real cuando $2-\omega^2=0$, en $\omega=\sqrt2\rads$. Allí $L=K/(-3\cdot2)=-K/6$.

\textbf{Baja frecuencia.} $\mathrm{Re}\,L\to-K(1+0{,}5)/2=-0{,}75K$ y $\mathrm{Im}\,L\to-\infty$ (tipo 1). El integrador cierra
la curva con un semicírculo infinito por la derecha.

\textbf{Conclusión.} Si $K/6<1$, el punto $-1$ queda afuera: $N=0$, $Z=0$ y el lazo es \textbf{estable}.
Si $K>6$, la curva rodea dos veces al $-1$ en sentido horario: $N=2$ y $Z=2$, \textbf{inestable}.
La ganancia crítica es $K_c=6$. Routh da lo mismo: $s^3+3s^2+2s+K$ es estable si $3\cdot2>K$.
\end{ejemplo}

\begin{figure}[ht]
\centering
\begin{tikzpicture}
\begin{axis}[width=0.6\linewidth,axis equal image,xmin=-3,xmax=1,ymin=-2.2,ymax=2.2,
  xlabel={Re $L(\jw)$},ylabel={Im $L(\jw)$},legend pos=outer north east]
\addplot[azul,thick] table[x=re,y=im]{datos/nyquist_k2.dat};
\addplot[rojo,thick] table[x=re,y=im]{datos/nyquist_k8.dat};
\addplot[azul,thick,dashed,forget plot] table[x=re,y=imneg]{datos/nyquist_k2.dat};
\addplot[rojo,thick,dashed,forget plot] table[x=re,y=imneg]{datos/nyquist_k8.dat};
\addplot[only marks,mark=x,mark size=3.5pt,very thick,black] coordinates{(-1,0)};
\legend{$K=2$ (estable),$K=8$ (inestable),$-1$}
\node[font=\scriptsize,azul,anchor=north] at (axis cs:-0.33,-0.05) {$-\tfrac13$};
\node[font=\scriptsize,rojo,anchor=north] at (axis cs:-1.45,-0.05) {$-\tfrac43$};
\end{axis}
\end{tikzpicture}
\caption{Nyquist del ejemplo~\ref{ej:nyq}. En línea llena $\omega>0$ y en línea discontinua $\omega<0$. Con $K=8$ el cruce cae a la izquierda de $-1$.}
\label{fig:nyq}
\end{figure}

% ---------------------------------------------------------------------
\section{Estabilidad relativa: márgenes de ganancia y de fase}
\label{sec:margenes}

Saber que el sistema es estable no alcanza: interesa \emph{cuán lejos} está de volverse inestable. Esa distancia
mide la robustez frente a errores de modelo y anticipa el amortiguamiento. Se definen dos frecuencias:
\begin{itemize}
\item $\wc$, \textbf{frecuencia de cruce de ganancia}: $|L(j\wc)|=1$ ($0\dB$).
\item $\wf$, \textbf{frecuencia de cruce de fase}: $\angle L(j\wf)=-180^\circ$.
\end{itemize}
\begin{formula}
\[
\MF=180^\circ+\angle L(j\wc),\qquad\qquad \MG=\frac{1}{|L(j\wf)|}\quad\Longleftrightarrow\quad \MG_{\mathrm{dB}}=-20\log|L(j\wf)|.
\]
\end{formula}
\begin{itemize}
\item El \textbf{MF} es el atraso de fase adicional que el lazo tolera en $\wc$ antes de llegar a la inestabilidad.
\item El \textbf{MG} es el factor por el que se puede multiplicar la ganancia antes de llegar a la inestabilidad.
      De ahí sale la \textbf{ganancia crítica}: $K_c=K\cdot10^{\MG_{\mathrm{dB}}/20}$.
\item En el diagrama polar, el MF es el ángulo entre el punto donde $L$ corta el círculo unitario y el semieje real negativo.
      $1/\MG$ es la distancia del cruce con el eje real al origen (figura~\ref{fig:margenes}).
\end{itemize}

\begin{figure}[ht]
\centering
\begin{tikzpicture}
\begin{groupplot}[group style={group size=1 by 2,vertical sep=0.3cm,xticklabels at=edge bottom},
  bode,width=0.56\linewidth,height=4.1cm,xmin=0.1,xmax=30]
\nextgroupplot[ylabel={$|L|$ [dB]},ymin=-60,ymax=30,ytick={-60,-40,-20,0,20}]
\addplot[exacta] table[x=w,y=mag]{datos/margenes.dat};
\draw[black] (axis cs:0.1,0)--(axis cs:30,0);
\draw[verde,dashed] (axis cs:0.749,-60)--(axis cs:0.749,30);
\draw[rojo,dashed] (axis cs:1.414,-60)--(axis cs:1.414,30);
\draw[rojo,thick,-{Stealth[length=4pt]}] (axis cs:1.414,0)--(axis cs:1.414,-9.54);
\node[font=\scriptsize,rojo,anchor=west] at (axis cs:1.5,-6) {$\MG=9{,}54\dB$};
\nextgroupplot[ylabel={$\angle L$ [$^\circ$]},xlabel={$\omega$ [rad/s]},ymin=-270,ymax=-90,ytick={-270,-225,-180,-135,-90}]
\addplot[exacta] table[x=w,y=fase]{datos/margenes.dat};
\draw[black] (axis cs:0.1,-180)--(axis cs:30,-180);
\draw[verde,dashed] (axis cs:0.749,-270)--(axis cs:0.749,-90);
\draw[rojo,dashed] (axis cs:1.414,-270)--(axis cs:1.414,-90);
\draw[verde,thick,{Stealth[length=4pt]}-] (axis cs:0.749,-147.4)--(axis cs:0.749,-180);
\node[font=\scriptsize,verde,anchor=east] at (axis cs:0.72,-160) {$\MF=32{,}6^\circ$};
\node[font=\scriptsize,verde,anchor=north east] at (axis cs:0.74,-232) {$\wc$};
\node[font=\scriptsize,rojo,anchor=north west] at (axis cs:1.45,-232) {$\wf$};
\end{groupplot}
\end{tikzpicture}\hfill
\begin{tikzpicture}
\begin{axis}[width=0.42\linewidth,height=0.42\linewidth,axis equal image,xmin=-1.6,xmax=0.6,ymin=-1.6,ymax=0.6,
  xlabel={Re},ylabel={Im}]
\addplot[gris,densely dotted] table[x=x,y=y]{datos/circulo.dat};
\addplot[azul,thick] table[x=re,y=im]{datos/nyquist_k2.dat};
\addplot[only marks,mark=x,mark size=3pt,very thick] coordinates{(-1,0)};
\draw[verde,thick] (axis cs:0,0)--(axis cs:-0.843,-0.539);
\draw[verde,thick,-{Stealth[length=4pt]}] (axis cs:-0.6,0) arc[start angle=180,end angle=212.6,radius=0.6];
\node[font=\scriptsize,verde] at (axis cs:-0.45,-0.16) {MF};
\draw[rojo,thick,{Stealth[length=4pt]}-{Stealth[length=4pt]}] (axis cs:-0.333,0.1)--(axis cs:0,0.1);
\node[font=\scriptsize,rojo,anchor=south] at (axis cs:-0.17,0.1) {$1/\MG$};
\end{axis}
\end{tikzpicture}
\caption{Márgenes de $L(s)=2/[s(s+1)(s+2)]$ sobre el Bode (izquierda) y sobre el diagrama polar, junto al círculo unitario (derecha).}
\label{fig:margenes}
\end{figure}

\subsection{Valores de diseño y margen de retardo}
Valores habituales de diseño: $\MF$ entre $30^\circ$ y $60^\circ$ (típicamente $45^\circ$--$50^\circ$) y $\MG\ge6\dB$ (típicamente $10$--$12\dB$).
Con $\MF<30^\circ$ la respuesta es muy oscilatoria. Con $\MF>70^\circ$ suele ser lenta.

Un retardo $\tau$ no cambia $|L|$ pero resta $\omega\tau$ radianes de fase. Si se le resta a $\wc$ todo el MF, el sistema
queda al borde de la inestabilidad. Por eso el \textbf{margen de retardo} es
\begin{formula}
\[
\tau_{\max}=\frac{\MF\ [\mathrm{rad}]}{\wc}=\frac{\MF[^\circ]}{57{,}3\,\wc}.
\]
\end{formula}
Para la figura~\ref{fig:margenes}: $\tau_{\max}=(32{,}6/57{,}3)/0{,}749=0{,}76\,$s.

\begin{cuidado}[Cuándo los márgenes engañan]
\begin{itemize}
\item \textbf{Varios cruces.} Si $|L|$ cruza $0\dB$ varias veces (por ejemplo, por una resonancia), hay varios MF y vale el menor.
\item \textbf{Lazo abierto inestable ($P\ne0$).} El MG puede salir negativo en dB aunque el sistema sea estable. Hay que usar Nyquist completo.
\item \textbf{Un margen solo no alcanza.} El ejemplo~\ref{ej:bode} tiene $\MF=106^\circ$ y $\MG=4{,}5\dB$. Una medida que resume ambos es la
      distancia mínima de la curva al punto $-1$: $\min_\omega|1+L(\jw)|=1/M_s$, donde $M_s$ es el pico de la sensibilidad. Se busca $M_s\lesssim2$ ($6\dB$).
\end{itemize}
\end{cuidado}

% ---------------------------------------------------------------------
\section{Respuesta en frecuencia de lazo cerrado}

\subsection{De $L$ a $T$}
Con realimentación unitaria, $T(\jw)=\dfrac{L(\jw)}{1+L(\jw)}$. Hay dos regiones sencillas:
\begin{itemize}
\item Donde $|L|\gg1$ (baja frecuencia), $T\approx1$: la salida sigue a la referencia.
\item Donde $|L|\ll1$ (alta frecuencia), $T\approx L$: el lazo cerrado filtra igual que el lazo abierto.
\end{itemize}
Entre ambas regiones, cerca de $\wc$, el comportamiento depende de la fase: si el MF es chico, $|1+L|$ es chico y $T$ presenta un pico.
Para caracterizar $T$ se usan:
\begin{itemize}
\item \textbf{Pico de resonancia} $M_r=\max|T(\jw)|$ y su frecuencia $\omega_r$. Un $M_r$ grande indica poco amortiguamiento.
      Valores buenos: $1{,}0<M_r<1{,}4$ ($0$--$3\dB$).
\item \textbf{Ancho de banda} $\omega_{BW}$: frecuencia donde $|T|$ cae $3\dB$ por debajo de su valor en $\omega=0$. Mide la rapidez: $t_r\cdot\omega_{BW}\approx$ cte.
      En la práctica, $\wc\lesssim\omega_{BW}\lesssim2\wc$.
\end{itemize}

\subsection{Carta de Nichols}
En el plano magnitud [dB] vs.\ fase [$^\circ$] se dibujan los \textbf{contornos M}, que son las curvas de $|T|=$ cte.
Al superponer $L(\jw)$, se lee $M_r$ como el contorno de mayor valor al que la curva queda tangente, y $\omega_{BW}$ donde la curva corta el contorno de $-3\dB$.
El punto crítico $-1$ es $(-180^\circ,\ 0\dB)$. El MF es la distancia horizontal del cruce de $0\dB$ a $-180^\circ$, y el MG es la distancia vertical
del cruce de $-180^\circ$ a $0\dB$. Su gran ventaja: \textbf{cambiar la ganancia desplaza la curva verticalmente sin deformarla}.

\begin{figure}[ht]
\centering
\begin{tikzpicture}
\begin{axis}[width=0.78\linewidth,height=6.6cm,xmin=-265,xmax=-90,ymin=-28,ymax=22,
  xtick={-270,-225,-180,-135,-90},xlabel={$\angle L$ [$^\circ$]},ylabel={$|L|$ [dB]},grid=major]
\addplot[gris!70,thin] table[x=fase,y=mag,empty line=jump]{datos/nichols_M.dat};
\addplot[only marks,mark=none,nodes near coords,point meta=explicit symbolic,
  every node near coord/.style={font=\tiny,gris,anchor=south}] table[x=fase,y=mag,meta=lab]{datos/nichols_lab.dat};
\addplot[azul,very thick] table[x=fase,y=mag]{datos/nichols_L.dat};
\addplot[only marks,mark=x,mark size=3.5pt,very thick,rojo] coordinates{(-180,0)};
\node[font=\scriptsize,azul,anchor=west] at (axis cs:-118,15) {$L(\jw)$, $\omega\uparrow$};
\end{axis}
\end{tikzpicture}
\caption{Carta de Nichols con $L=2/[s(s+1)(s+2)]$. Los contornos grises son $|T|$ constante en dB. La curva queda tangente a un contorno
intermedio entre los de $3$ y $6\dB$: $M_r=5{,}3\dB$ en $\omega_r=0{,}82\rads$, coherente con su $\MF=32{,}6^\circ$.}
\label{fig:nichols}
\end{figure}

\subsection{Correlación tiempo--frecuencia}
Para el lazo de segundo orden estándar
\[
L(s)=\frac{\omega_n^2}{s(s+2\zeta\omega_n)}\qquad\Longrightarrow\qquad T(s)=\frac{\omega_n^2}{s^2+2\zeta\omega_n s+\omega_n^2}
\]
se obtienen relaciones exactas:
\begin{formula}
\[
\MF=\arctan\frac{2\zeta}{\sqrt{\sqrt{1+4\zeta^4}-2\zeta^2}}\approx100\,\zeta\quad(\zeta<0{,}7),
\]
\[
M_p=e^{-\pi\zeta/\sqrt{1-\zeta^2}},\qquad M_r=\frac{1}{2\zeta\sqrt{1-\zeta^2}},\qquad t_s(2\pct{})\approx\frac{4}{\zeta\omega_n},
\]
\[
\wc=\omega_n\sqrt{\sqrt{1+4\zeta^4}-2\zeta^2},\qquad
\omega_{BW}=\omega_n\sqrt{1-2\zeta^2+\sqrt{4\zeta^4-4\zeta^2+2}}.
\]
\end{formula}
La regla práctica es \textbf{$\zeta\approx\MF/100$} (con el MF en grados). A partir de ella se estima el sobrepico.
Si el sistema tiene polos dominantes, las relaciones valen aproximadamente para sistemas de orden mayor.

\begin{figure}[ht]
\centering
\begin{tikzpicture}
\begin{groupplot}[group style={group size=2 by 1,horizontal sep=1.6cm},width=0.48\linewidth,height=4.8cm,xmin=0,xmax=1,xlabel={$\zeta$}]
\nextgroupplot[ylabel={MF [$^\circ$]},ymin=0,ymax=80,legend pos=south east]
\addplot[azul,thick] table[x=z,y=mf]{datos/correlacion.dat};
\addplot[rojo,dashed,thick] table[x=z,y=aprox]{datos/correlacion.dat};
\legend{Exacta,$100\zeta$}
\nextgroupplot[ylabel={$M_p$ [\%]},ymin=0,ymax=100]
\addplot[azul,thick] table[x=z,y=mp]{datos/correlacion.dat};
\end{groupplot}
\end{tikzpicture}
\caption{Margen de fase y sobrepico en función de $\zeta$ para el segundo orden estándar.}
\label{fig:correl}
\end{figure}

\begin{center}
\small
\begin{tabular}{@{}lcccccc@{}}
\toprule
$\zeta$ & 0,2 & 0,3 & 0,4 & 0,5 & 0,6 & 0,7\\
\midrule
MF & $22{,}6^\circ$ & $33{,}3^\circ$ & $43{,}1^\circ$ & $51{,}8^\circ$ & $59{,}2^\circ$ & $65{,}2^\circ$\\
$M_p$ & 52,7\pct{} & 37,2\pct{} & 25,4\pct{} & 16,3\pct{} & 9,5\pct{} & 4,6\pct{}\\
$M_r$ [dB] & 8,14 & 4,85 & 2,70 & 1,25 & 0,35 & 0,00\\
\bottomrule
\end{tabular}
\end{center}

% ---------------------------------------------------------------------
\section{Ejercicios de la Parte I}
\label{sec:ejI}

\ejercicio{Trazado asintótico}\label{e:bode}
Sea $G(s)=\dfrac{500\,(s+2)}{s\,(s+10)\,(s+20)}$.
(a)~Llevarla a forma de Bode y armar la tabla de quiebres.
(b)~Trazar la magnitud asintótica y dar su valor en $\omega=1$, 2, 10 y 20~rad/s.
(c)~Estimar $\wc$ con las asíntotas y comparar con el valor exacto, $16{,}6\rads$.
(d)~Calcular la fase en $\omega=10\rads$.
(e)~Dar $K_v$ y el error ante rampa unitaria.

\ejercicio{Segundo orden}\label{e:2do}
Para $G(s)=\dfrac{100}{s^2+4s+100}$, calcular $\zeta$, $\omega_n$, $M_r$ (en dB), $\omega_r$, el error de la asíntota en
$\omega_n$, la fase en $\omega_n$ y el ancho de banda.

\ejercicio{Identificación}\label{e:ident}
Un Bode medido de un sistema de fase mínima tiene asíntotas de $-20\dBdec$ hasta $\omega=2$, de $0\dBdec$ entre 2 y
20, y de $-40\dBdec$ de ahí en adelante. La prolongación del tramo inicial corta $0\dB$ en $\omega=10\rads$. En
$\omega=20$, la curva real está $6\dB$ \emph{por debajo} de la asíntota. Encontrar $G(s)$.
¿Cómo cambiaría la respuesta si la curva real estuviera $8\dB$ \emph{por encima}?

\ejercicio{Fase no mínima y retardo}\label{e:fnm}
Comparar $G_1(s)=\dfrac{s+1}{s+5}$, $G_2(s)=\dfrac{1-s}{s+5}$ y $G_3(s)=\dfrac{e^{-0{,}2s}}{s+5}$.
(a)~¿Qué tienen en común sus magnitudes?
(b)~Calcular la fase de $G_1$ y de $G_2$ en $\omega=1$, en $\omega=5$ y para $\omega\to\infty$.
(c)~¿Cuánto atrasa el retardo en $\omega=5$?
(d)~¿Cuáles son de fase mínima?

\ejercicio{Nyquist, sistema tipo 0}\label{e:nyq0}
Para $L(s)=\dfrac{K}{(s+1)^3}$ con $K>0$:
(a)~esbozar el diagrama polar indicando el punto inicial, el cruce con el eje real y la dirección de llegada;
(b)~encontrar el rango de $K$ que da estabilidad;
(c)~calcular el MG para $K=2$.

\ejercicio{Nyquist con lazo abierto inestable}\label{e:nyqP}
Sea $L(s)=\dfrac{K\,(s+3)}{s\,(s-1)}$ con $K>0$.
(a)~¿Cuánto vale $P$?
(b)~Encontrar el cruce con el eje real.
(c)~Usar $Z=N+P$ para obtener el rango de $K$ estable y verificar con el polinomio característico.
(d)~¿Qué significa que MATLAB informe $\MG=-6\dB$ para $K=2$?

\ejercicio{Márgenes y retardo}\label{e:marg}
Para $L(s)=\dfrac{10}{s\,(1+s/5)\,(1+s/50)}$ calcular $\wf$, MG, $K_c$, $\wc$, MF y el margen de retardo.

\ejercicio{Correlación}\label{e:corr}
En un ensayo de lazo cerrado se midió $M_r=3\dB$ en $\omega_r=5\rads$. Suponiendo polos dominantes de segundo orden,
estimar $\zeta$, $\omega_n$, $M_p$, $t_s$ (2\pct{}) y el MF del lazo abierto.

% =====================================================================
\parte{II}{Compensación en frecuencia}
% =====================================================================

\section{Planteo del problema}

\subsection{Compensación serie y especificaciones}
Se agrega un compensador $G_c(s)$ en serie con la planta $G(s)$, de modo que el lazo abierto es $L=G_cG$:
\begin{center}
\begin{tikzpicture}[>=Stealth,node distance=1.4cm,
  blk/.style={draw,thick,minimum height=8mm,minimum width=14mm,fill=azul!6},
  sum/.style={draw,thick,circle,inner sep=2pt}]
\node (r) {$R(s)$};
\node[sum,right=1cm of r] (s) {};
\node[blk,right=1cm of s] (gc) {$G_c(s)$};
\node[blk,right=1.2cm of gc] (g) {$G(s)$};
\node[right=1.2cm of g] (y) {$Y(s)$};
\draw[->,thick] (r)--(s) node[pos=0.85,above]{\scriptsize$+$};
\draw[->,thick] (s)--node[above]{\small$E$}(gc);
\draw[->,thick] (gc)--node[above]{\small$U$}(g);
\draw[->,thick] (g)--coordinate[pos=0.4](a)(y);
\draw[->,thick] (a)--++(0,-1.1)-|(s) node[pos=0.95,right]{\scriptsize$-$};
\end{tikzpicture}
\end{center}
Las especificaciones típicas en frecuencia, y lo que cada una controla, son:
\begin{center}
\small
\begin{tabular}{@{}lll@{}}
\toprule
Especificación & Región de $|L|$ & Se relaciona con\\
\midrule
$K_p$, $K_v$, $K_a$ & baja frecuencia & error en régimen, rechazo de perturbaciones lentas\\
MF, MG ($M_r$) & cerca de $\wc$ & amortiguamiento, sobrepico, robustez\\
$\wc$, $\omega_{BW}$ & cerca de $\wc$ & velocidad ($t_r$, $t_s$)\\
Atenuación en alta frecuencia & $\omega\gg\wc$ & ruido de medición, dinámica no modelada\\
\bottomrule
\end{tabular}
\end{center}

\subsection{La forma ideal del lazo}
\begin{idea}[Moldeo del lazo (\emph{loop shaping})]
Un buen $|L(\jw)|$ tiene:
\begin{itemize}
\item \textbf{ganancia alta en baja frecuencia} para que el error sea chico;
\item \textbf{cruce de $0\dB$ con pendiente de alrededor de $-20\dBdec$}, mantenida un tramo de casi una década. Por la relación de Bode,
      esto da una fase cercana a $-90^\circ$ más el atraso de los polos vecinos, es decir, un MF sano. Si cruza con $-40\dBdec$,
      la fase está cerca de $-180^\circ$ y el MF es casi nulo;
\item \textbf{caída rápida en alta frecuencia} para no amplificar ruido.
\end{itemize}
Compensar es \textbf{deformar} $|L|$ y $\angle L$ para que se parezcan a esta forma.
\end{idea}

\begin{figure}[ht]
\centering
\begin{tikzpicture}[>=Stealth,x=1.25cm,y=0.045cm]
\fill[verde!10] (0,0) rectangle (3,60);
\fill[rojo!8] (6.2,-80) rectangle (9.4,0);
\draw[->] (0,0)--(9.6,0) node[right]{\small$\log\omega$};
\draw[->] (0,-80)--(0,62) node[above]{\small$|L|$ [dB]};
\draw[azul,very thick] (0.2,55)--(3,30)--(4.2,8)--(5.4,-8)--(6.2,-20)--(9.2,-75);
\node[verde!50!black,font=\small,align=center] at (1.6,15) {ganancia alta:\\ error chico};
\node[rojo,font=\small,align=center] at (7.9,-12) {ganancia baja:\\ ruido rechazado};
\draw[dashed,gris] (4.8,62)--(4.8,-80) node[below]{\small$\wc$};
\node[azul,font=\small,anchor=west] at (5.0,24) {pendiente $\approx-20\dBdec$ en el cruce};
\node[gris,font=\scriptsize,anchor=north] at (1.6,-5) {(especificación de $K_v$)};
\node[gris,font=\scriptsize,anchor=north] at (7.9,-30) {(especificación de ruido)};
\end{tikzpicture}
\caption{Esquema de la forma deseada de $|L(\jw)|$. La curva debe esquivar las dos zonas sombreadas y cruzar $0\dB$ suavemente.}
\label{fig:shape}
\end{figure}

\subsection{Lo que no puede hacer la ganancia sola}
Un compensador puramente proporcional, $G_c=K$, sube o baja toda la curva de magnitud sin tocar la fase. Al subir $K$
para cumplir el error, $\wc$ se corre a la derecha, donde la fase es peor, y el MF cae. Al bajar $K$ para recuperar MF,
se pierde precisión. \textbf{Para cumplir ambas especificaciones hace falta modificar la forma de la curva, y eso
requiere polos y ceros.} Ese es el trabajo de las redes de adelanto, de atraso y de atraso-adelanto.

\begin{receta}[Primer paso, común a los tres diseños]
Fijar la ganancia $K$ que cumple la especificación de error (por ejemplo, $K_v=\lim_{s\to0}sKG(s)$) y estudiar el Bode de
$L_0=KG$, \emph{el sistema ya ajustado en ganancia pero sin compensar}. Todos los compensadores de este apunte tienen
ganancia unitaria en $\omega=0$, así que no alteran el $K$ fijado.
\end{receta}

% ---------------------------------------------------------------------
\section{Compensador de adelanto de fase}

\subsection{Estructura y diagrama de Bode}
\begin{formula}
\[
G_c(s)=K\,\frac{Ts+1}{\alpha Ts+1}=\frac{K}{\alpha}\cdot\frac{s+1/T}{s+1/(\alpha T)},\qquad 0<\alpha<1.
\]
\end{formula}
El cero ($1/T$) está más cerca del origen que el polo ($1/(\alpha T)$). Entre ambos, la red se comporta como un derivador:
\textbf{sube la magnitud} $20\log(1/\alpha)\dB$ en total y \textbf{aporta fase positiva}, con un máximo $\phi_m$ a mitad de camino
en escala logarítmica (figura~\ref{fig:leadnorm}).

\begin{figure}[ht]
\centering
\begin{tikzpicture}
\begin{groupplot}[group style={group size=2 by 1,horizontal sep=1.6cm},
  bode,width=0.48\linewidth,height=4.8cm,xmin=0.01,xmax=1000,xlabel={$\omega T$}]
\nextgroupplot[ylabel={$|\frac{\jw T+1}{\jw\alpha T+1}|$ [dB]},ymin=-1,ymax=21,ytick={0,6,12,20},legend pos=north west]
\addplot[verde,thick] table[x=x,y=m05]{datos/adelanto_norm.dat};
\addplot[naranja,thick] table[x=x,y=m025]{datos/adelanto_norm.dat};
\addplot[rojo,thick] table[x=x,y=m01]{datos/adelanto_norm.dat};
\legend{$\alpha=0{,}5$,$\alpha=0{,}25$,$\alpha=0{,}1$}
\nextgroupplot[ylabel={Fase [$^\circ$]},ymin=0,ymax=60,ytick={0,19.5,36.9,54.9}]
\addplot[verde,thick] table[x=x,y=f05]{datos/adelanto_norm.dat};
\addplot[naranja,thick] table[x=x,y=f025]{datos/adelanto_norm.dat};
\addplot[rojo,thick] table[x=x,y=f01]{datos/adelanto_norm.dat};
\draw[gris,densely dotted] (axis cs:3.162,0)--(axis cs:3.162,54.9);
\node[font=\scriptsize,gris,anchor=south] at (axis cs:3.162,55.5) {$\omega_m T=1/\sqrt{\alpha}$};
\end{groupplot}
\end{tikzpicture}
\caption{Red de adelanto normalizada. La fase máxima se da en la media geométrica del cero y el polo.}
\label{fig:leadnorm}
\end{figure}

\subsection{Fórmulas de diseño}
La fase de la red es $\phi(\omega)=\arctan(\omega T)-\arctan(\alpha\omega T)$. Derivando e igualando a cero:
\begin{formula}
\[
\omega_m=\frac{1}{T\sqrt\alpha}=\sqrt{\frac1T\cdot\frac{1}{\alpha T}},\qquad
\sin\phi_m=\frac{1-\alpha}{1+\alpha}\ \Longleftrightarrow\ \alpha=\frac{1-\sin\phi_m}{1+\sin\phi_m},
\]
\[
\left|\frac{j\omega_mT+1}{j\omega_m\alpha T+1}\right|=\frac{1}{\sqrt\alpha}\quad\Longleftrightarrow\quad+10\log\frac1\alpha\ \mathrm{dB}.
\]
\end{formula}
Para la última igualdad: $\big|\frac{j/\sqrt\alpha+1}{j\sqrt\alpha+1}\big|^2=\frac{1+1/\alpha}{1+\alpha}=\frac1\alpha$.
En $\omega_m$ la red aporta \emph{la mitad} (en dB) de su ganancia total de alta frecuencia.
\begin{center}
\small
\begin{tabular}{@{}l*{8}{c}@{}}
\toprule
$\phi_m$ & $20^\circ$ & $30^\circ$ & $35^\circ$ & $40^\circ$ & $45^\circ$ & $50^\circ$ & $55^\circ$ & $60^\circ$\\
$\alpha$ & 0,490 & 0,333 & 0,271 & 0,217 & 0,172 & 0,132 & 0,0994 & 0,0718\\
$10\log(1/\alpha)$ [dB] & 3,1 & 4,8 & 5,7 & 6,6 & 7,7 & 8,8 & 10,0 & 11,4\\
\bottomrule
\end{tabular}
\end{center}

\subsection{Qué hace y qué no hace}
\begin{itemize}
\item \textbf{Mejora:} aumenta el MF, aumenta $\wc$ y el ancho de banda (respuesta \textbf{más rápida}) y reduce el sobrepico.
\item \textbf{Costo:} aumenta la ganancia de alta frecuencia en $1/\alpha$, por lo que amplifica el ruido de medición y exige
      mayor esfuerzo de control (posible saturación del actuador).
\item \textbf{Limitaciones:} en la práctica $\phi_m\lesssim60^\circ$ por etapa ($\alpha\gtrsim0{,}07$). Si hace falta más, se ponen
      dos redes en cascada. Si la fase de la planta cae muy rápido cerca del cruce (polos cercanos, retardo), el adelanto
      corre $\wc$ hacia donde la fase es todavía peor y la ganancia de fase se pierde. En ese caso conviene el atraso.
\end{itemize}

\subsection{Procedimiento de diseño}
\begin{receta}[Diseño de adelanto por Bode]
\begin{enumerate}[label=\textbf{\arabic*.}]
\item Elegir $K$ para cumplir el error en régimen. Trazar el Bode de $L_0=KG$ y obtener $\MF_0$ y $\omega_{c0}$.
\item Fase a agregar: $\phi_m=\MF_{\text{req}}-\MF_0+\delta$, con $\delta\approx5^\circ$--$12^\circ$. El $\delta$ compensa que el cruce
      se corre a la derecha, donde $L_0$ tiene menos fase. Usar un $\delta$ mayor cuanto más rápido cae la fase de $L_0$.
\item Calcular $\alpha=\dfrac{1-\sin\phi_m}{1+\sin\phi_m}$.
\item \textbf{Nuevo cruce:} buscar la frecuencia en la que $|L_0(\jw)|_{\mathrm{dB}}=-10\log(1/\alpha)$. Esa frecuencia es $\omega_m$
      y será el nuevo $\wc$, porque ahí la red suma justo $+10\log(1/\alpha)$ y la magnitud total da $0\dB$.
      Así se aprovecha el máximo de fase exactamente en el cruce.
\item Ubicar la red: $1/T=\omega_m\sqrt\alpha$ (cero) y $1/(\alpha T)=\omega_m/\sqrt\alpha$ (polo).
\item Armar $G_c(s)=K\dfrac{Ts+1}{\alpha Ts+1}$ y \textbf{verificar} MF y MG con el Bode compensado.
      Si no se cumple, volver al paso 2 con un $\delta$ mayor.
\end{enumerate}
\end{receta}

\begin{ejemplo}{Diseño de adelanto (con iteración)}
\label{ej:lead}
Planta: $G(s)=\dfrac{1}{s\,(s+1)\,(0{,}05s+1)}$. Especificaciones: $K_v=10\,\mathrm{s^{-1}}$, $\MF\ge45^\circ$ y $\MG\ge10\dB$.

\textbf{1. Ganancia.} $K_v=\lim sKG=K$, así que $K=10$. Para $L_0=10G$ se obtiene $\omega_{c0}=3{,}07\rads$,
$\MF_0=9{,}4^\circ$ y $\MG_0=6{,}4\dB$. El sistema es estable pero muy oscilatorio ($M_p\approx77\pct{}$).

\textbf{2--4. Primer intento, con $\delta=5^\circ$.} $\phi_m=45-9{,}4+5=40{,}6^\circ$, de donde $\alpha=0{,}211$ y $10\log(1/\alpha)=6{,}8\dB$.
$|L_0|=-6{,}8\dB$ en $\omega_m=4{,}55\rads$. Al verificar, $\MF=40{,}2^\circ$: \textbf{no cumple}.
La causa es que entre 3,07 y 4,55~rad/s la fase de $L_0$ cayó casi $10^\circ$ (el polo en 1 y el polo en 20 siguen restando),
más que los $5^\circ$ de reserva.

\textbf{2--4. Segundo intento, con $\delta=12^\circ$.} $\phi_m=47{,}6^\circ$, de donde $\alpha=0{,}15$ y $10\log(1/\alpha)=8{,}2\dB$.
$|L_0(\jw)|=-8{,}2\dB$ en $\omega_m=4{,}96\rads$. Verificación de fase: $\angle L_0(j4{,}96)=-90-\arctan4{,}96-\arctan0{,}248=-182{,}5^\circ$,
y al sumarle la red, $-182{,}5+47{,}6=-134{,}9^\circ$. Resulta $\MF\approx45{,}1^\circ$. \checkmark

\textbf{5. Ubicación.} Cero en $1/T=4{,}96\sqrt{0{,}15}=1{,}92\rads$ y polo en $1/(\alpha T)=4{,}96/\sqrt{0{,}15}=12{,}8\rads$.
\[
\boxed{G_c(s)=10\,\frac{0{,}521\,s+1}{0{,}0781\,s+1}=66{,}7\,\frac{s+1{,}92}{s+12{,}8}}
\]
\textbf{6. Verificación (simulación).} $\MF=45{,}2^\circ$, $\MG=14{,}9\dB$ y $\wc=4{,}95\rads$. $K_v=10$ se mantiene. El sobrepico baja de
77\pct{} a 27\pct{}, el tiempo de establecimiento de 15,3~s a 1,0~s, y el ancho de banda sube de 4,8 a 8,9~rad/s
(figura~\ref{fig:leadej}).
\end{ejemplo}

\begin{figure}[ht]
\centering
\begin{tikzpicture}
\begin{groupplot}[group style={group size=1 by 2,vertical sep=0.3cm,xticklabels at=edge bottom},
  bode,width=0.56\linewidth,height=4.2cm,xmin=0.1,xmax=300]
\nextgroupplot[ylabel={Magnitud [dB]},ymin=-80,ymax=45,legend pos=south west,legend columns=1]
\addplot[antes] table[x=w,y=m0]{datos/adelanto_ej.dat};
\addplot[despues] table[x=w,y=m1]{datos/adelanto_ej.dat};
\addplot[comp] table[x=w,y=mc]{datos/adelanto_ej.dat};
\legend{$L_0=10G$,$L=G_cG$,$G_c/K$}
\draw[black] (axis cs:0.1,0)--(axis cs:300,0);
\nextgroupplot[ylabel={Fase [$^\circ$]},xlabel={$\omega$ [rad/s]},ymin=-270,ymax=60,ytick={-270,-180,-135,-90,0,45}]
\addplot[antes] table[x=w,y=f0]{datos/adelanto_ej.dat};
\addplot[despues] table[x=w,y=f1]{datos/adelanto_ej.dat};
\addplot[comp] table[x=w,y=fc]{datos/adelanto_ej.dat};
\draw[black] (axis cs:0.1,-180)--(axis cs:300,-180);
\draw[verde,dashed] (axis cs:4.95,-270)--(axis cs:4.95,60);
\end{groupplot}
\end{tikzpicture}\hfill
\begin{tikzpicture}
\begin{axis}[width=0.41\linewidth,height=6.6cm,xmin=0,xmax=8,ymin=0,ymax=1.85,xlabel={$t$ [s]},ylabel={$y(t)$},legend pos=north east]
\addplot[antes] table[x=t,y=y0]{datos/adelanto_ej_step.dat};
\addplot[despues] table[x=t,y=y1]{datos/adelanto_ej_step.dat};
\legend{Sin comp.,Adelanto}
\end{axis}
\end{tikzpicture}
\caption{Ejemplo~\ref{ej:lead}. Izquierda: el adelanto suma fase justo en el nuevo cruce (línea verde). Derecha: respuesta al escalón de lazo cerrado.}
\label{fig:leadej}
\end{figure}

% ---------------------------------------------------------------------
\section{Compensador de atraso de fase}

\subsection{Estructura y diagrama de Bode}
\begin{formula}
\[
G_c(s)=K\,\frac{Ts+1}{\beta Ts+1}=\frac{K}{\beta}\cdot\frac{s+1/T}{s+1/(\beta T)},\qquad \beta>1.
\]
\end{formula}
Ahora el \textbf{polo} ($1/(\beta T)$) está más cerca del origen que el cero ($1/T$). En baja frecuencia la red vale 1 (0 dB),
y en alta frecuencia vale $1/\beta$, es decir, \textbf{atenúa} $20\log\beta\dB$. Su fase es negativa, que es justamente lo que \emph{no} se quiere,
así que se la ubica \textbf{lejos, a la izquierda de $\wc$}, donde ese atraso ya no molesta.

\begin{figure}[ht]
\centering
\begin{tikzpicture}
\begin{groupplot}[group style={group size=2 by 1,horizontal sep=1.6cm},
  bode,width=0.48\linewidth,height=4.4cm,xmin=0.001,xmax=100,xlabel={$\omega T$}]
\nextgroupplot[ylabel={Magnitud [dB]},ymin=-22,ymax=2,ytick={-20,-10,0}]
\addplot[azul,thick] table[x=x,y=mag]{datos/atraso_norm.dat};
\draw[gris,densely dotted] (axis cs:0.1,-22)--(axis cs:0.1,2) (axis cs:1,-22)--(axis cs:1,2);
\node[font=\scriptsize,anchor=north west] at (axis cs:0.11,1.5) {$\frac{1}{\beta T}$};
\node[font=\scriptsize,anchor=north west] at (axis cs:1.1,1.5) {$\frac{1}{T}$};
\node[font=\scriptsize,azul,anchor=south] at (axis cs:15,-19.5) {$-20\log\beta$};
\nextgroupplot[ylabel={Fase [$^\circ$]},ymin=-60,ymax=2,ytick={-54.9,-40,-20,-5.1,0}]
\addplot[azul,thick] table[x=x,y=fase]{datos/atraso_norm.dat};
\draw[rojo,dashed] (axis cs:10,-60)--(axis cs:10,2);
\node[font=\scriptsize,rojo,anchor=north east,align=right] at (axis cs:9,-30) {en $\omega=10/T$\\ solo $-5^\circ$};
\end{groupplot}
\end{tikzpicture}
\caption{Red de atraso normalizada con $\beta=10$. Una década por encima del cero ya atenúa casi los $20\dB$ completos y atrasa apenas $5^\circ$.}
\label{fig:lagnorm}
\end{figure}

\subsection{La idea}
El atraso no se usa por su fase: \textbf{se usa por su atenuación}. Se busca la frecuencia donde $L_0$ tiene la fase que
daría el MF pedido. En general esa frecuencia está a la izquierda del cruce actual, y $L_0$ tiene ahí una magnitud de
$A\dB>0$. Con $\beta=10^{A/20}$, la red baja la curva justo $A\dB$ en esa zona, y esa frecuencia pasa a ser el nuevo cruce.
Como el cero se ubica una década antes, el atraso que la red aporta en el cruce es de solo $\approx5^\circ$, que se reservan
de antemano. La ganancia de baja frecuencia (el $K_v$) no cambia.

\subsection{Procedimiento de diseño}
\begin{receta}[Diseño de atraso por Bode]
\begin{enumerate}[label=\textbf{\arabic*.}]
\item Elegir $K$ para cumplir el error. Trazar el Bode de $L_0=KG$.
\item Buscar la frecuencia $\wc'$ en la que $\angle L_0(j\wc')=-180^\circ+\MF_{\text{req}}+\delta$, con $\delta\approx5^\circ$--$12^\circ$
      (típicamente $6^\circ$). Será el nuevo cruce. Se la puede redondear hacia abajo, que es el lado seguro.
\item Leer $A=|L_0(j\wc')|_{\mathrm{dB}}$ y tomar $\beta=10^{A/20}$.
\item Cero una década antes del nuevo cruce: $1/T=\wc'/10$ (entre $\wc'/10$ y $\wc'/5$).
\item Polo en $1/(\beta T)$.
\item Armar $G_c(s)=K\dfrac{Ts+1}{\beta Ts+1}$ y verificar MF y MG.
\end{enumerate}
\end{receta}

\begin{ejemplo}{Diseño de atraso}
\label{ej:lag}
Planta: $G(s)=\dfrac{1}{s\,(s+2)\,(s+5)}$. Especificaciones: $K_v=10\,\mathrm{s^{-1}}$, $\MF\ge45^\circ$ y $\MG\ge10\dB$.

\textbf{1. Ganancia.} $K_v=K/(2\cdot5)=10$, así que $K=100$. Con $L_0=100G$ se obtiene $\MF_0=-8{,}9^\circ$ y $\MG_0=-3{,}1\dB$:
\textbf{inestable} ($K_c=70$ según Routh). Un adelanto tendría que aportar más de $60^\circ$ en una zona donde la fase cae rápido, así que se elige atraso.

\textbf{2. Nuevo cruce.} Se busca $\angle L_0=-180+45+6=-129^\circ$:
$-90^\circ-\arctan(\omega/2)-\arctan(\omega/5)=-129^\circ$, que da $\omega\approx1{,}03$. Se redondea a $\wc'=1\rads$.
Ahí $\angle L_0(j1)=-90-26{,}6-11{,}3=-127{,}9^\circ$: hay $52{,}1^\circ$ disponibles antes de la red.

\textbf{3. Atenuación.} $|L_0(j1)|=\dfrac{100}{1\cdot\sqrt5\cdot\sqrt{26}}=8{,}77$, o sea $A=18{,}9\dB$. Se toma $\beta=9$, redondeando hacia arriba.

\textbf{4--5. Ubicación.} Cero en $1/T=0{,}1\rads$ ($T=10$\,s) y polo en $1/(\beta T)=0{,}0111\rads$.
\[
\boxed{G_c(s)=100\,\frac{10\,s+1}{90\,s+1}=11{,}1\,\frac{s+0{,}1}{s+0{,}0111}}
\]
\textbf{6. Verificación.} $\MF=47{,}5^\circ$, $\MG=15{,}4\dB$ y $\wc=0{,}98\rads$, con $K_v=10$. \checkmark

\textbf{Comparación con la ganancia sola.} Para tener $\MF=45^\circ$ bajando solo la ganancia, haría falta $K=14{,}7$, o sea $K_v=1{,}47$,
con un error de rampa de 0,68 contra 0,10 del atraso. Las dos respuestas al escalón se parecen, pero el atraso sigue la rampa
casi siete veces mejor (figura~\ref{fig:lagej}).
\end{ejemplo}

\begin{figure}[ht]
\centering
\begin{tikzpicture}
\begin{groupplot}[group style={group size=1 by 2,vertical sep=0.3cm,xticklabels at=edge bottom},
  bode,width=0.92\linewidth,height=4cm,xmin=0.001,xmax=100]
\nextgroupplot[ylabel={Magnitud [dB]},ymin=-80,ymax=80,ytick={-80,-40,0,40,80},legend pos=south west]
\addplot[antes] table[x=w,y=m0]{datos/atraso_ej.dat};
\addplot[despues] table[x=w,y=m1]{datos/atraso_ej.dat};
\addplot[comp] table[x=w,y=mc]{datos/atraso_ej.dat};
\legend{$L_0=100G$,$L=G_cG$,$G_c/K$}
\draw[black] (axis cs:0.001,0)--(axis cs:100,0);
\draw[gris,densely dotted] (axis cs:0.0111,-80)--(axis cs:0.0111,80) (axis cs:0.1,-80)--(axis cs:0.1,80);
\nextgroupplot[ylabel={Fase [$^\circ$]},xlabel={$\omega$ [rad/s]},ymin=-270,ymax=0,ytick={-270,-180,-135,-90,0}]
\addplot[antes] table[x=w,y=f0]{datos/atraso_ej.dat};
\addplot[despues] table[x=w,y=f1]{datos/atraso_ej.dat};
\addplot[comp] table[x=w,y=fc]{datos/atraso_ej.dat};
\draw[black] (axis cs:0.001,-180)--(axis cs:100,-180);
\draw[verde,dashed] (axis cs:0.98,-270)--(axis cs:0.98,0);
\end{groupplot}
\end{tikzpicture}

\begin{tikzpicture}
\begin{groupplot}[group style={group size=2 by 1,horizontal sep=1.6cm},width=0.49\linewidth,height=4.4cm,xmin=0,xmax=25,xlabel={$t$ [s]}]
\nextgroupplot[ylabel={$y(t)$ (escalón)},ymin=0,ymax=1.4,legend pos=south east]
\addplot[despues] table[x=t,y=y1]{datos/atraso_ej_step.dat};
\addplot[naranja,thick,dashed] table[x=t,y=yg]{datos/atraso_ej_step.dat};
\legend{Atraso ($K_v=10$),Solo ganancia ($K_v=1{,}47$)}
\nextgroupplot[ylabel={$e(t)$ (rampa)},ymin=0,ymax=1]
\addplot[despues] table[x=t,y=e1]{datos/atraso_ej_rampa.dat};
\addplot[naranja,thick,dashed] table[x=t,y=eg]{datos/atraso_ej_rampa.dat};
\end{groupplot}
\end{tikzpicture}
\caption{Ejemplo~\ref{ej:lag}. Arriba: el atraso baja la curva $19\dB$ sin tocar la baja frecuencia. El lazo cruza en $\approx1\rads$,
donde todavía hay fase. Abajo: escalón y error ante rampa contra el ajuste por ganancia sola.}
\label{fig:lagej}
\end{figure}

\subsection{Efectos y precauciones}
\begin{itemize}
\item \textbf{Mejora} el error en régimen sin sacrificar el MF y \textbf{atenúa el ruido} de alta frecuencia.
\item \textbf{Reduce} $\wc$ y el ancho de banda: la respuesta es \textbf{más lenta}.
\item El par polo-cero cerca del origen (un \textbf{dipolo}) deja un polo lento de lazo cerrado con residuo chico. Eso produce
      una \emph{cola} que tarda en asentarse, y por eso el $t_s$ del ejemplo es de unos 15~s aunque el sobrepico sea moderado.
      Cuanto más cerca del origen está el dipolo, más larga es la cola.
\item Si $\beta$ resulta enorme (más de 20 o 30) o el nuevo cruce queda demasiado bajo para la velocidad pedida, conviene pasar al atraso-adelanto.
\end{itemize}

% ---------------------------------------------------------------------
\section{Compensador de atraso-adelanto}

\subsection{Estructura}
Es el producto de una red de atraso y una de adelanto:
\begin{formula}
\[
G_c(s)=K\,\underbrace{\frac{T_1s+1}{\alpha T_1s+1}}_{\text{adelanto}}\cdot\underbrace{\frac{T_2s+1}{\beta T_2s+1}}_{\text{atraso}},\qquad
\alpha<1<\beta,\qquad \frac{1}{\beta T_2}<\frac1{T_2}<\frac1{T_1}<\frac1{\alpha T_1}.
\]
\end{formula}
Se reparte el trabajo: la parte de \textbf{adelanto} aporta fase en el cruce (MF y velocidad), y la de \textbf{atraso} aporta la
atenuación que hace falta para que el cruce caiga donde se quiere, sin tocar la ganancia de baja frecuencia (error).
Se usa cuando se piden a la vez precisión, buen MF y una velocidad que el atraso solo no logra,
o cuando el adelanto solo necesitaría una fase imposible.

\subsection{Procedimiento de diseño}
\begin{receta}[Diseño de atraso-adelanto por Bode]
\begin{enumerate}[label=\textbf{\arabic*.}]
\item Elegir $K$ para cumplir el error. Trazar el Bode de $L_0=KG$.
\item Elegir el nuevo cruce $\wc'$. Una elección clásica es $\wc'\approx\wf$ de $L_0$ (o una especificación de ancho de banda),
      con lo que la fase de $L_0$ ahí es $\approx-180^\circ$.
\item \textbf{Adelanto:} $\phi_m=\MF_{\text{req}}-\big(180^\circ+\angle L_0(j\wc')\big)+\delta$, con $\delta\approx6^\circ$, que es el atraso que
      luego mete la otra red. Calcular $\alpha$ y centrar la red en $\wc'$: $1/T_1=\wc'\sqrt\alpha$ y $1/(\alpha T_1)=\wc'/\sqrt\alpha$.
\item \textbf{Atraso:} calcular $A=|L_0(j\wc')|_{\mathrm{dB}}+10\log(1/\alpha)$, que es lo que sobra en $\wc'$ después del adelanto,
      y tomar $\beta=10^{A/20}$. Cero en $1/T_2=\wc'/10$ y polo en $1/(\beta T_2)$.
\item Verificar MF, MG y la respuesta temporal.
\end{enumerate}
\end{receta}

\begin{ejemplo}{Diseño de atraso-adelanto}
\label{ej:laglead}
Planta: $G(s)=\dfrac{1}{s\,(s+1)\,(s+4)}$. Especificaciones: $K_v=10\,\mathrm{s^{-1}}$, $\MF\ge50^\circ$ y $\MG\ge10\dB$.

\textbf{1.} $K_v=K/4=10$, así que $K=40$. $L_0=40G$ es \textbf{inestable}: $\MF_0=-15^\circ$ y $\MG_0=-6\dB$. Su $\wf$ es 2~rad/s,
porque $\arctan\omega+\arctan(\omega/4)=90^\circ$ cuando $\omega^2=4$.

\textbf{2.} Se elige $\wc'=2\rads$. Ahí $\angle L_0=-180^\circ$ y $|L_0(j2)|=\dfrac{40}{2\sqrt5\sqrt{20}}=2$, es decir $6{,}02\dB$.

\textbf{3. Adelanto.} $\phi_m=50-0+6=56^\circ$, de donde $\alpha=\dfrac{1-\sin56^\circ}{1+\sin56^\circ}=0{,}0935$ y $10\log(1/\alpha)=10{,}3\dB$.
Cero en $2\sqrt{0{,}0935}=0{,}611$ y polo en $2/\sqrt{0{,}0935}=6{,}54$.

\textbf{4. Atraso.} $A=6{,}02+10{,}3=16{,}3\dB$, de donde $\beta=6{,}54$. Cero en $2/10=0{,}2$ y polo en $0{,}2/6{,}54=0{,}0306$.
\[
\boxed{G_c(s)=40\,\frac{(s/0{,}611+1)\,(s/0{,}2+1)}{(s/6{,}54+1)\,(s/0{,}0306+1)}=65{,}5\,\frac{(s+0{,}611)(s+0{,}2)}{(s+6{,}54)(s+0{,}0306)}}
\]
\textbf{5. Verificación.} $\MF=51{,}0^\circ$, $\MG=13{,}3\dB$ y $\wc=2{,}01\rads$, con $K_v=10$. $M_p=18\pct{}$ y ancho de banda de 3,6~rad/s. \checkmark

\textbf{¿Por qué no uno solo?} Con atraso solo, el cruce caería donde $\angle L_0=-124^\circ$, cerca de 0,5~rad/s: cuatro veces más lento.
Con adelanto solo, harían falta más de $70^\circ$ de fase.
\end{ejemplo}

\begin{figure}[ht]
\centering
\begin{tikzpicture}
\begin{groupplot}[group style={group size=1 by 2,vertical sep=0.3cm,xticklabels at=edge bottom},
  bode,width=0.56\linewidth,height=4.1cm,xmin=0.001,xmax=100]
\nextgroupplot[ylabel={Magnitud [dB]},ymin=-80,ymax=80,ytick={-80,-40,0,40,80},legend pos=south west]
\addplot[antes] table[x=w,y=m0]{datos/ad_at_ej.dat};
\addplot[despues] table[x=w,y=m1]{datos/ad_at_ej.dat};
\addplot[comp] table[x=w,y=mc]{datos/ad_at_ej.dat};
\legend{$L_0$,$L$,$G_c/K$}
\draw[black] (axis cs:0.001,0)--(axis cs:100,0);
\nextgroupplot[ylabel={Fase [$^\circ$]},xlabel={$\omega$ [rad/s]},ymin=-270,ymax=60,ytick={-270,-180,-130,-90,0,45}]
\addplot[antes] table[x=w,y=f0]{datos/ad_at_ej.dat};
\addplot[despues] table[x=w,y=f1]{datos/ad_at_ej.dat};
\addplot[comp] table[x=w,y=fc]{datos/ad_at_ej.dat};
\draw[black] (axis cs:0.001,-180)--(axis cs:100,-180);
\draw[verde,dashed] (axis cs:2.01,-270)--(axis cs:2.01,60);
\end{groupplot}
\end{tikzpicture}\hfill
\begin{tikzpicture}
\begin{axis}[width=0.41\linewidth,height=6.4cm,xmin=0,xmax=20,ymin=0,ymax=1.3,xlabel={$t$ [s]},ylabel={$y(t)$}]
\addplot[despues] table[x=t,y=y1]{datos/ad_at_ej_step.dat};
\end{axis}
\end{tikzpicture}
\caption{Ejemplo~\ref{ej:laglead}. La red combinada atenúa en la zona media (atraso) y suma fase alrededor de 2~rad/s (adelanto).}
\label{fig:lagleadej}
\end{figure}

% ---------------------------------------------------------------------
\section{Comparación, relación con el PID y criterios de elección}

\subsection{Misma planta, dos soluciones}
Para $G(s)=1/[s(s+1)]$ con $K_v=10$ y $\MF\ge45^\circ$, el sistema solo con ganancia tiene $\MF=18^\circ$ y $M_p=60\pct{}$.
Se lo puede compensar de las dos maneras:
\begin{itemize}
\item \textbf{Adelanto:} $\alpha=0{,}307$, cero en 2,32 y polo en 7,57. Resulta $\wc=4{,}19$, $\MF=45{,}5^\circ$, $M_p=27\pct{}$ y $t_s=1{,}7$\,s.
\item \textbf{Atraso:} $\beta=9{,}6$, cero en 0,081 y polo en 0,0084. Resulta $\wc=0{,}81$, $\MF=45{,}8^\circ$, $M_p=26\pct{}$ y $t_s=15{,}5$\,s.
\end{itemize}
Ambos cumplen las especificaciones, con sobrepicos casi iguales porque los MF son iguales. \textbf{La diferencia es la velocidad}:
el adelanto cruza cinco veces más arriba (figura~\ref{fig:comp}).

\begin{figure}[ht]
\centering
\begin{tikzpicture}
\begin{groupplot}[group style={group size=2 by 1,horizontal sep=1.6cm},width=0.49\linewidth,height=5cm]
\nextgroupplot[xmode=log,minor x tick num=8,xmin=0.001,xmax=100,ymin=-60,ymax=80,xlabel={$\omega$ [rad/s]},ylabel={$|L|$ [dB]},legend pos=south west]
\addplot[antes] table[x=w,y=m0]{datos/comparacion.dat};
\addplot[rojo,thick] table[x=w,y=mad]{datos/comparacion.dat};
\addplot[azul,thick] table[x=w,y=mat]{datos/comparacion.dat};
\legend{Solo $K$,Adelanto,Atraso}
\draw[black] (axis cs:0.001,0)--(axis cs:100,0);
\nextgroupplot[xmin=0,xmax=15,ymin=0,ymax=1.7,xlabel={$t$ [s]},ylabel={$y(t)$}]
\addplot[antes] table[x=t,y=y0]{datos/comparacion_step.dat};
\addplot[rojo,thick] table[x=t,y=yad]{datos/comparacion_step.dat};
\addplot[azul,thick] table[x=t,y=yat]{datos/comparacion_step.dat};
\end{groupplot}
\end{tikzpicture}
\caption{Adelanto contra atraso sobre $G=1/[s(s+1)]$ con $K_v=10$. El adelanto sube $\wc$ y el atraso lo baja.}
\label{fig:comp}
\end{figure}

\begin{center}
\small
\begin{tabularx}{\linewidth}{@{}l>{\raggedright\arraybackslash}X>{\raggedright\arraybackslash}X>{\raggedright\arraybackslash}X@{}}
\toprule
 & Adelanto & Atraso & Atraso-adelanto\\
\midrule
Mecanismo & suma fase cerca de $\wc$ & atenúa en media y alta frecuencia & ambos\\
$\wc$, ancho de banda & aumentan (más rápido) & disminuyen (más lento) & según el diseño\\
MF & aumenta & aumenta (cruza donde hay fase) & aumenta\\
Error en régimen & lo fija $K$ & lo fija $K$ (se puede subir sin perder MF) & lo fija $K$\\
Ruido de alta frecuencia & lo amplifica ($\times1/\alpha$) & lo atenúa ($\times1/\beta$) & intermedio\\
Esfuerzo de control & mayor & menor & intermedio\\
Respuesta & rápida, mayor $u(t)$ inicial & cola lenta por el dipolo & cola lenta leve\\
Equivalente PID & PD con filtro & PI aproximado & PID aproximado\\
\bottomrule
\end{tabularx}
\end{center}

\subsection{Relación con PD, PI y PID}
\begin{itemize}
\item \textbf{PD:} $G_c=K_p(1+T_ds)$ es un adelanto con $\alpha\to0$: suma hasta $+90^\circ$ pero su ganancia crece sin límite
      ($+20\dBdec$). En la práctica se le agrega un polo, $\frac{1+T_ds}{1+T_ds/N}$, y queda un adelanto con $\alpha=1/N$.
\item \textbf{PI:} $G_c=K_p\Big(1+\dfrac{1}{T_is}\Big)=K_p\dfrac{T_is+1}{T_is}$ es un atraso con $\beta\to\infty$: el polo va al origen.
      Sube el tipo en uno (error nulo al escalón si la planta era tipo 0). Se diseña como el atraso: cero en $1/T_i\approx\wc/10$,
      lo que cuesta unos $5{,}7^\circ$ de fase en el cruce.
\item \textbf{PID:} combina ambos: $G_c=K_p\dfrac{T_iT_ds^2+T_is+1}{T_is}$. Los dos ceros cumplen el papel de los ceros del atraso-adelanto.
\end{itemize}

\subsection{Cómo elegir}
\begin{receta}[Árbol de decisión rápido]
\begin{enumerate}[label=\textbf{\arabic*.}]
\item Con $K$ fijado por el error: ¿el MF ya alcanza? Entonces no hace falta compensador.
\item ¿Falta poco MF (menos de $50^\circ$--$60^\circ$), la fase de $L_0$ no cae abruptamente cerca de $\wc$ y se quiere más velocidad (o no molesta)?
      Usar \textbf{adelanto}.
\item ¿La velocidad actual sobra, hay ruido o la fase cae muy rápido cerca del cruce? Usar \textbf{atraso}.
\item ¿Hacen falta precisión, MF y velocidad a la vez, o el adelanto solo pide más de $60^\circ$? Usar \textbf{atraso-adelanto} (o dos adelantos).
\end{enumerate}
\end{receta}

\begin{cuidado}[Errores frecuentes]
\begin{itemize}
\item Diseñar sobre $G$ en lugar de sobre $L_0=KG$. \textbf{Primero se fija la ganancia y después se compensa}.
\item Ubicar el máximo de fase del adelanto en el cruce \emph{viejo}: el cruce se mueve a $\omega_m$.
\item Olvidar la reserva $\delta$, o no volver a verificar. Siempre hay que recalcular el MF con el lazo completo.
\item En el atraso, poner el cero demasiado cerca de $\wc'$: resta mucha fase. Una década antes es lo seguro.
\item Confundir $\alpha$ con $\beta$, o el cero con el polo. En el adelanto el \emph{cero} está más cerca del origen; en el atraso, el \emph{polo}.
\item Leer $K_v$ sobre la curva real en lugar de la asíntota.
\end{itemize}
\end{cuidado}

% ---------------------------------------------------------------------
\section{Ejercicios de la Parte II}
\label{sec:ejII}

\ejercicio{Adelanto}\label{p:lead}
Para $G(s)=\dfrac{1}{s\,(s+4)}$, diseñar un adelanto que dé $K_v=20\,\mathrm{s^{-1}}$ y $\MF\ge50^\circ$. Probar primero con $\delta=5^\circ$.

\ejercicio{Atraso en un sistema tipo 0}\label{p:lag}
Para $G(s)=\dfrac{1}{(s+1)(s+2)(s+5)}$, se pide error al escalón $\le2\pct{}$ y $\MF\ge40^\circ$.
(a)~Calcular $K$ y mostrar que el lazo resulta inestable.
(b)~Diseñar un atraso con $\delta=6^\circ$, redondeando el nuevo cruce a un valor entero.
(c)~Verificar que el error sigue siendo del 2\pct{}.

\ejercicio{Atraso-adelanto}\label{p:laglead}
Para $G(s)=\dfrac{1}{s\,(s+1)\,(s+2)}$, se pide $K_v=10\,\mathrm{s^{-1}}$, $\MF\ge50^\circ$ y $\MG\ge10\dB$.
Diseñar un atraso-adelanto tomando $\wc'=\wf$ de $L_0$ y $\delta=6^\circ$.

\ejercicio{Elegir el compensador}\label{p:elegir}
Para cada caso, indicar si conviene adelanto, atraso o atraso-adelanto, y por qué.
(a)~Cumple $K_v$ con $\MF=30^\circ$, la fase de $L_0$ cae suavemente y se quiere bajar el tiempo de establecimiento.
(b)~Cumple MF y velocidad, pero el error de rampa es 5 veces mayor que el permitido.
(c)~Con la $K$ requerida, $\MF=-20^\circ$, y el sensor tiene mucho ruido de alta frecuencia.
(d)~Con la $K$ requerida, $\MF=-15^\circ$, y se pide además que $\wc$ no baje.

\ejercicio{Retardo}\label{p:retardo}
El lazo $L(s)=\dfrac{K\,e^{-\tau s}}{s\,(s+1)}$ se ajustó con $K=2$ y $\tau=0$.
(a)~Calcular $\wc$, el MF y el retardo máximo admisible.
(b)~Si en la práctica hay un retardo $\tau=0{,}2$\,s, ¿cuál es el nuevo MF?
(c)~¿Qué $K$ hace falta para tener $\MF=30^\circ$ con ese retardo? ¿Cuánto vale $K_v$?

\ejercicio{PI como atraso}\label{p:pi}
La planta $G(s)=\dfrac{10}{(s+1)(0{,}1s+1)}$ con control proporcional unitario tiene un error al escalón de $1/11$.
Diseñar un PI $G_c=K_p\frac{T_is+1}{T_is}$ con $K_p=1$ que anule ese error y no reste más de $6^\circ$ de MF.
Calcular el MF antes y después.

\ejercicio{Problema integrador: la aeronave de la UT05}\label{p:avion}
Se retoma el caso de estudio de la UT05, el cabeceo de una aeronave con fuselaje flexible:
\[
L(s)=K\,G_aG_\theta G_f,\qquad G_a=\frac{50}{s+50},\quad G_\theta=\frac{18(s+1)}{s(s^2+4{,}2s+9)},\quad G_f=\frac{400}{s^2+s+400}.
\]
$G_f$ es el modo elástico del fuselaje: $\omega_f=20\rads$ y $\zeta_f=0{,}025$.
Con $K=1$ se obtiene $K_v=2$, $\MF=45{,}8^\circ$, $\MG=14{,}3\dB$ y un pico elástico de $-1{,}5\dB$ en $|L|$.
Se quiere \textbf{duplicar la precisión ante rampa} ($K_v=4$), con $\MF\ge45^\circ$ y con el pico elástico de $|L|$
\textbf{por lo menos $6\dB$ por debajo de $0\dB$}, para que ese modo no cruce $0\dB$ aunque $\zeta_f$ sea menor que la nominal.
\begin{enumerate}[label=(\alph*)]
\item Con $K=2$ y sin compensar: ¿cuánto vale el pico elástico? ¿Por qué ya no alcanza con mirar el MF que informa \texttt{margin}?
\item Intentar un adelanto para $\MF=45^\circ$ con $\delta=10^\circ$. ¿Qué le pasa al modo elástico?
\item Diseñar un atraso. ¿Qué especificación manda: el MF o el pico elástico?
\item Verificar y comparar con el diseño nominal ($K=1$).
\end{enumerate}

% =====================================================================
\section{Soluciones}
\label{sec:soluciones}
\subsection*{Parte I}

\solucion{e:bode}{Trazado asintótico}
(a)~$G=\dfrac{500\cdot2}{10\cdot20}\cdot\dfrac{1+s/2}{s(1+s/10)(1+s/20)}=\dfrac{5(1+s/2)}{s(1+s/10)(1+s/20)}$, con $K_B=5$ y tipo 1.
Quiebres: cero en 2 ($+20$), polo en 10 ($-20$) y polo en 20 ($-20$). Pendiente final: $-40\dBdec$ ($n-m=2$).

(b)~$\omega<2$: $20\log(5/\omega)$, que en $\omega=1$ vale $\mathbf{13{,}98\dB}$ y en $\omega=2$ vale $\mathbf{7{,}96\dB}$.
Entre 2 y 10 la curva es plana en $7{,}96\dB$, así que en $\omega=10$ vale $\mathbf{7{,}96\dB}$.
Entre 10 y 20 baja a $-20\dBdec$: $7{,}96-20\log2=\mathbf{1{,}94\dB}$ en $\omega=20$. De ahí en adelante, $-40\dBdec$.
Los valores exactos son 14,9, 10,8, 4,1 y $-2{,}0\dB$: la asíntota falla en unos 3~dB cerca de cada quiebre.

(c)~La asíntota de $-40$ cruza $0\dB$ en $20\cdot10^{1{,}94/40}=22{,}4\rads$. El exacto es 16,6: la asíntota sobreestima porque en 10 y en 20
la curva real ya está unos $3\dB$ por debajo. Con varios quiebres cercanos, hay que usar la expresión exacta.

(d)~$\angle G(j10)=-90+\arctan5-\arctan1-\arctan0{,}5=-90+78{,}7-45-26{,}6=\rta{-82{,}9^\circ}$.

(e)~$K_v=K_B=5\,\mathrm{s^{-1}}$ y $e_{ss}=1/K_v=0{,}2$.

\solucion{e:2do}{Segundo orden}
$\omega_n=10\rads$ y $2\zeta\omega_n=4$, así que $\zeta=0{,}2$.
$M_r=\dfrac{1}{2(0{,}2)\sqrt{1-0{,}04}}=2{,}55$, es decir $\rta{8{,}14\dB}$, en $\omega_r=10\sqrt{1-0{,}08}=\rta{9{,}59\rads}$.
En $\omega_n$: $|G|=1/(2\zeta)=2{,}5$, o sea $7{,}96\dB$ por encima de la asíntota. La fase es $-90^\circ$.
$\omega_{BW}=10\sqrt{1-0{,}08+\sqrt{0{,}0064-0{,}16+2}}=10\sqrt{0{,}92+1{,}359}=\rta{15{,}1\rads}$.

\solucion{e:ident}{Identificación}
Tipo 1 (pendiente inicial $-20$), con la asíntota en $0\dB$ en $\omega=10$, así que $K_B=10$.
De 2 a 20 la pendiente sube 20: cero en 2. En 20 baja 40: dos polos.
Una diferencia de $-6\dB=-20\log2$ en el quiebre indica un polo real doble ($\zeta=1$, porque $-20\log(2\cdot1)=-6$).
\[
G(s)=\frac{10\,(1+s/2)}{s\,(1+s/20)^2}=\rta{\frac{2000\,(s+2)}{s\,(s+20)^2}}
\]
Si la curva estuviera $8\dB$ \emph{por encima}: $-20\log(2\zeta)=8$, de donde $\zeta=0{,}2$, y el denominador tendría $(s^2+8s+400)$ en lugar de $(s+20)^2$.

\solucion{e:fnm}{Fase no mínima y retardo}
(a)~Como $|1-\jw|=|1+\jw|$, $G_1$ y $G_2$ tienen magnitudes \textbf{idénticas}: $\sqrt{1+\omega^2}/\sqrt{25+\omega^2}$.
Como $|e^{-0{,}2\jw}|=1$, $G_3$ tiene la misma magnitud que $1/(s+5)$. En los tres casos el Bode de magnitud no muestra nada raro.

(b)~$\angle G_1=\arctan\omega-\arctan(\omega/5)$ y $\angle G_2=-\arctan\omega-\arctan(\omega/5)$.
\begin{center}
\small
\begin{tabular}{@{}lccc@{}}
\toprule
 & $\omega=1$ & $\omega=5$ & $\omega\to\infty$\\
\midrule
$G_1$ & $45-11{,}3=33{,}7^\circ$ & $78{,}7-45=33{,}7^\circ$ & $0^\circ$\\
$G_2$ & $-45-11{,}3=-56{,}3^\circ$ & $-78{,}7-45=-123{,}7^\circ$ & $-180^\circ$\\
\bottomrule
\end{tabular}
\end{center}
(c)~$-0{,}2\cdot5=-1\,$rad, es decir $\rta{-57{,}3^\circ}$ sin cambio de magnitud.
(d)~Solo $G_1$. $G_2$ tiene un cero en $s=+1$ y $G_3$ un retardo. Ambos agregan atraso de fase sin que la magnitud lo delate.

\solucion{e:nyq0}{Nyquist, sistema tipo 0}
(a)~Arranca en $L(0)=K$ sobre el eje real positivo. La fase es $-3\arctan\omega$: cruza el eje imaginario negativo ($-90^\circ$)
en $\omega=\tan30^\circ=1/\sqrt3$ y el eje real negativo ($-180^\circ$) en $\omega=\tan60^\circ=\sqrt3$.
Allí $|L|=K/(\sqrt{1+3})^3=K/8$. La curva llega al origen tangente a $-270^\circ$ (la dirección $+j$).

(b)~$P=0$. Para no rodear $-1$ hace falta $K/8<1$: $\rta{0<K<8}$.

(c)~Con $K=2$ el cruce está en $-1/4$: $\MG=4$, es decir $\rta{12{,}0\dB}$. (El MF es $67{,}6^\circ$ en $\wc=0{,}77\rads$.)

\solucion{e:nyqP}{Nyquist con lazo abierto inestable}
(a)~El polo en $s=+1$ está en el SPD: $P=1$. El polo en $s=0$ se esquiva con la indentación.

(b)~$L(\jw)=\dfrac{K(3+\jw)}{-\omega^2-\jw}=\dfrac{-K\big[4\omega+j(\omega^2-3)\big]}{\omega(\omega^2+1)}$. La parte imaginaria se anula en $\omega=\sqrt3$,
y ahí $L=-\dfrac{4\sqrt3K}{\sqrt3\cdot4}=-K$. Para $\omega\to0^+$, $\mathrm{Re}\to-4K$ e $\mathrm{Im}\to+\infty$.

(c)~La curva para $\omega>0$ baja desde $(-4K,+j\infty)$, cruza en $-K$ y llega al origen por el semiplano inferior.
Con el espejo y el semicírculo infinito (horario, por la izquierda), el lazo chico entre $-K$ y el origen se recorre en sentido
\textbf{antihorario}. Si $K>1$, el $-1$ queda adentro del lazo chico: $N=-1$ y $Z=N+P=0$, estable.
Si $K<1$, el $-1$ queda en el lazo grande, recorrido en sentido horario: $N=+1$ y $Z=2$, inestable.
Polinomio característico: $s^2+(K-1)s+3K$, estable si y solo si $\rta{K>1}$. \checkmark

\begin{center}
\begin{tikzpicture}
\begin{axis}[width=0.62\linewidth,axis equal image,xmin=-9,xmax=1,ymin=-4.2,ymax=4.2,
  xlabel={Re},ylabel={Im},grid=major,legend pos=outer north east,legend style={font=\scriptsize}]
\addplot[azul,thick,postaction={decorate},decoration={markings,
  mark=at position 0.25 with {\arrow{Stealth[length=5pt]}},mark=at position 0.75 with {\arrow{Stealth[length=5pt]}}}]
  table[x=re,y=im]{datos/nyquist_inest.dat};
\addplot[azul,thick,dashed,postaction={decorate},decoration={markings,
  mark=at position 0.25 with {\arrowreversed{Stealth[length=5pt]}},mark=at position 0.75 with {\arrowreversed{Stealth[length=5pt]}}}]
  table[x=re,y=imneg]{datos/nyquist_inest.dat};
\addplot[only marks,mark=x,mark size=3.5pt,very thick,rojo] coordinates{(-1,0)};
\legend{$\omega>0$,$\omega<0$,$-1$}
\node[font=\scriptsize,anchor=north] at (axis cs:-2,-0.15) {$-K$};
\node[font=\scriptsize,gris,align=center] at (axis cs:-5.6,0) {cierre por $\infty$\\(horario, por la izq.)};
\end{axis}
\end{tikzpicture}
\end{center}
\vspace{-4pt}
{\small Nyquist para $K=2$: el lazo chico entre $-2$ y $0$ se recorre en sentido antihorario y encierra al $-1$.}

(d)~Con $K=2$ el cruce está en $-2$. Para llegar a la inestabilidad hay que \emph{reducir} la ganancia a la mitad, hasta $K=1$.
El MG es entonces un margen de \emph{reducción} de ganancia: $20\log(1/2)=-6\dB$. No indica inestabilidad, lo cual confirma que con $P\ne0$ los márgenes se interpretan con Nyquist.

\solucion{e:marg}{Márgenes y retardo}
$\angle L=-90-\arctan(\omega/5)-\arctan(\omega/50)=-180$, de donde $\frac{\omega}{5}\cdot\frac{\omega}{50}=1$ y $\wf=\sqrt{250}=\rta{15{,}8\rads}$.
$|L(j\wf)|=\dfrac{10}{15{,}81\cdot\sqrt{11}\cdot\sqrt{1{,}1}}=\dfrac{10}{55}$, así que $\MG=5{,}5$, es decir $\rta{14{,}8\dB}$, y $K_c=10\cdot5{,}5=\rta{55}$.
Routh sobre $s^3+55s^2+250s+250K$ da $55\cdot250>250K$, o sea $K<55$. \checkmark
Cruce de ganancia (numérico): $\wc=\rta{6{,}22\rads}$, con $\angle L=-90-51{,}2-7{,}1=-148{,}3^\circ$, así que $\MF=\rta{31{,}7^\circ}$.
Margen de retardo: $\tau_{\max}=(31{,}7/57{,}3)/6{,}22=\rta{0{,}089\,\mathrm{s}}$.

\solucion{e:corr}{Correlación}
$M_r=10^{3/20}=1{,}413$. De $\frac{1}{2\zeta\sqrt{1-\zeta^2}}=1{,}413$ sale $\zeta^2(1-\zeta^2)=0{,}1253$, cuya raíz útil es $\zeta^2=0{,}1469$, o sea $\rta{\zeta=0{,}383}$.
$\omega_n=\omega_r/\sqrt{1-2\zeta^2}=5/0{,}840=\rta{5{,}95\rads}$.
$M_p=e^{-\pi(0{,}383)/0{,}924}=\rta{27\pct{}}$ y $t_s\approx4/(0{,}383\cdot5{,}95)=\rta{1{,}75\,\mathrm{s}}$.
$\MF=\arctan\dfrac{0{,}767}{\sqrt{1{,}0423-0{,}2938}}=\rta{41{,}5^\circ}$. La regla $100\zeta$ da $38^\circ$.

\subsection*{Parte II}

\solucion{p:lead}{Adelanto}
$K_v=K/4=20$, así que $K=80$. $L_0=80/[s(s+4)]$: $\omega_{c0}=8{,}51\rads$ y $\MF_0=90-\arctan(8{,}51/4)=25{,}2^\circ$.

\emph{Con $\delta=5^\circ$:} $\phi_m=29{,}8^\circ$, de donde $\alpha=0{,}336$ y $4{,}7\dB$. $|L_0|=-4{,}7\dB$ en $\omega_m=11{,}4$. MF obtenido: $49{,}1^\circ$. \textbf{No alcanza por poco.}

\emph{Con $\delta=7^\circ$:} $\phi_m=31{,}8^\circ$, de donde $\alpha=0{,}31$ y $5{,}1\dB$. $\omega_m=11{,}66\rads$.
El cero queda en $11{,}66\sqrt{0{,}31}=6{,}49$ y el polo en $11{,}66/\sqrt{0{,}31}=20{,}96$.
\[
G_c(s)=80\,\frac{s/6{,}49+1}{s/20{,}96+1}=\rta{258\,\frac{s+6{,}49}{s+20{,}96}}
\]
Resulta $\MF=50{,}8^\circ$ y $\wc=11{,}7\rads$. El MG es infinito porque la fase nunca llega a $-180^\circ$. $M_p\approx20\pct{}$.

\solucion{p:lag}{Atraso en un sistema tipo 0}
(a)~$K_p=K/10$ y $e_{ss}=1/(1+K_p)\le0{,}02$, así que $K_p=49$ y $K=490$.
$L_0=490G$ tiene $\MF=-32{,}3^\circ$ y $\MG=-11{,}8\dB$: \textbf{inestable}. Routh da $K_c=126$.

(b)~Fase objetivo: $-180+40+6=-134^\circ$, que ocurre en $\omega\approx2{,}2$. Se redondea a $\wc'=2\rads$, donde $\angle L_0=-63{,}4-45-21{,}8=-130{,}2^\circ$.
$|L_0(j2)|=\dfrac{490}{\sqrt5\sqrt8\sqrt{29}}=14{,}4$, o sea $23{,}2\dB$, y por lo tanto $\beta=14{,}4$. Cero en $0{,}2$ y polo en $0{,}2/14{,}4=0{,}0139$.
\[
G_c(s)=490\,\frac{5s+1}{72s+1}=\rta{34{,}0\,\frac{s+0{,}2}{s+0{,}0139}}
\]
Verificación: $\MF=44{,}2^\circ$, $\MG=10{,}6\dB$ y $\wc=2{,}0\rads$.

(c)~$G_c(0)=490$, así que $K_p$ sigue siendo 49 y el error sigue siendo del 2\pct{}. El atraso no toca la ganancia estática.

\solucion{p:laglead}{Atraso-adelanto}
$K=20$. $\wf=\sqrt2=1{,}414\rads$, donde $|L_0|=20/(\sqrt2\cdot\sqrt3\cdot\sqrt6)=3{,}33$, es decir $10{,}5\dB$.
Adelanto: $\phi_m=50+6=56^\circ$, así que $\alpha=0{,}0935$ ($+10{,}3\dB$). Cero en $1{,}414\sqrt{0{,}0935}=0{,}432$ y polo en $4{,}63$.
Atraso: $A=10{,}5+10{,}3=20{,}8\dB$, así que $\beta=10{,}9$. Cero en $0{,}141$ y polo en $0{,}0130$.
\[
G_c(s)=20\,\frac{(s/0{,}432+1)(s/0{,}141+1)}{(s/4{,}63+1)(s/0{,}013+1)}
\]
Verificación: $\MF=50{,}6^\circ$, $\MG=12{,}5\dB$, $\wc=1{,}42\rads$ y $M_p\approx18\pct{}$. \checkmark

\solucion{p:elegir}{Elegir el compensador}
(a)~\textbf{Adelanto}: falta poca fase, la fase no cae abruptamente y además se quiere más velocidad. El adelanto da ambas cosas.
(b)~\textbf{Atraso}, o un PI: hay que multiplicar por 5 la ganancia de baja frecuencia sin mover el cruce. Con $\beta=5$ y el cero una década antes de $\wc$,
el MF cae unos $5^\circ$ y la velocidad casi no cambia.
(c)~\textbf{Atraso}: hay que bajar mucho la magnitud para recuperar fase, y el ruido aconseja no levantar la alta frecuencia.
(d)~\textbf{Atraso-adelanto}: con $\MF=-15^\circ$ el adelanto solo pediría más de $60^\circ$, y el atraso solo bajaría $\wc$.

\solucion{p:retardo}{Retardo}
(a)~$\omega\sqrt{1+\omega^2}=2$ lleva a $\omega^4+\omega^2-4=0$, de donde $\omega^2=1{,}562$ y $\wc=\rta{1{,}25\rads}$.
$\MF=90-\arctan1{,}25=\rta{38{,}7^\circ}$ y $\tau_{\max}=0{,}675/1{,}25=\rta{0{,}54\,\mathrm{s}}$.

(b)~El retardo resta $0{,}2\cdot1{,}25=0{,}25\,$rad, es decir $14{,}3^\circ$: el MF baja a $\rta{24{,}3^\circ}$. El retardo no cambia $\wc$.

(c)~Se pide $90-\arctan\wc-11{,}46\,\wc=30$, con la fase del retardo en grados. Resolviendo numéricamente, $\wc=1{,}09\rads$, y de $|L|=1$ sale
$K=\wc\sqrt{1+\wc^2}=\rta{1{,}62}$. Entonces $K_v=1{,}62$: \textbf{el retardo obliga a resignar precisión}, salvo que se compense.

\solucion{p:pi}{PI como atraso}
Con P unitario, $|L|=1$ cuando $(1+\omega^2)(1+0{,}01\omega^2)=100$, es decir $\omega^4+101\omega^2-9900=0$. De ahí $\omega^2=61{,}1$ y $\wc=7{,}82\rads$.
$\MF=180-\arctan7{,}82-\arctan0{,}782=180-82{,}7-38{,}0=59{,}3^\circ$.
El PI aporta $-\arctan\frac{1}{\wc T_i}$ en el cruce. Con $T_i=10/\wc=1{,}28$\,s el aporte es $-\arctan0{,}1=-5{,}7^\circ$, y la magnitud
apenas sube ($\times1{,}005$), así que el cruce casi no se mueve.
\[
G_c(s)=\frac{1{,}28\,s+1}{1{,}28\,s}\qquad\Rightarrow\qquad \MF\approx59{,}3-5{,}7=\rta{53{,}6^\circ},\quad e_{ss}=0.
\]
La simulación da $\MF=53{,}5^\circ$. El PI hizo el sistema tipo 1 a cambio de menos de $6^\circ$ de fase.

\solucion{p:avion}{Problema integrador: la aeronave de la UT05}
(a)~$K=2$ sube toda la curva $6{,}02\dB$ y el pico elástico pasa de $-1{,}5$ a $\rta{+4{,}5\dB}$. Ahora $|L|$ cruza $0\dB$
\textbf{tres veces}: en 6,27, en 19,2 y en 20,6~rad/s. \texttt{margin} informa el MF del primer cruce ($23{,}8^\circ$), pero en los otros dos
la fase está muy lejos de $-180^\circ$ (\texttt{allmargin} da ``MF'' de $-43^\circ$ y $-153^\circ$). El lazo cerrado es estable (polos elásticos en
$-0{,}78\pm j19{,}1$, con $\zeta\approx0{,}04$), pero el modo elástico queda estabilizado \emph{solo por la fase}: si $\zeta_f$ o $\omega_f$ cambian un poco,
puede desestabilizarse. Además, $\MF=23{,}8^\circ$ tampoco cumple.

(b)~$\phi_m=45-23{,}8+10=31{,}2^\circ$, de donde $\alpha=0{,}318$, $\omega_m=8{,}85\rads$, cero en 4,99 y polo en 15,7. En 20~rad/s la red ya da casi toda su
ganancia de alta frecuencia: $+8{,}1\dB$. El pico elástico sube a $+12{,}6\dB$ y el lazo cerrado resulta \textbf{inestable}:
aparecen polos en $+0{,}61\pm j17{,}4$, es decir, el fuselaje oscila con amplitud creciente. \textbf{El adelanto levanta la ganancia justo donde vive la resonancia.}

(c)~Hay dos condiciones sobre la atenuación de la red:
\begin{itemize}
\item \emph{Por MF:} $\angle L_0=-180+45+6=-129^\circ$ en $\omega=3{,}83$, donde $|L_0|=7{,}1\dB$. Alcanza con $\beta=2{,}26$, pero el pico quedaría en
      $4{,}5-7{,}1=-2{,}6\dB$: \textbf{no cumple}.
\item \emph{Por el pico:} hace falta atenuar $4{,}5+6=10{,}5\dB$ en 20~rad/s, o sea $\beta\ge10^{10{,}5/20}=3{,}35$.
\end{itemize}
\textbf{Manda el pico elástico.} Con $\beta\ge3{,}35$, el nuevo cruce tiene que estar donde $|L_0|\ge10{,}5\dB$, es decir en $\omega\le2{,}75$.
Se elige $\wc'=2{,}5\rads$, donde $|L_0|=11{,}2\dB$ y $\angle L_0=-100{,}4^\circ$. Entonces $\beta=3{,}62$, el cero va en $0{,}25$ ($T=4$\,s) y el polo en $0{,}069$:
\[
G_c(s)=2\,\frac{4s+1}{14{,}5s+1}=\rta{0{,}552\,\frac{s+0{,}25}{s+0{,}069}}
\]
(d)~Resultados verificados por simulación:
\begin{center}
\small
\begin{tabular}{@{}lccccccc@{}}
\toprule
Diseño & $K_v$ & MF & MG & $\wc$ [rad/s] & Pico elástico & $M_p$ & $t_s$ [s]\\
\midrule
Nominal $K=1$ & 2 & $45{,}8^\circ$ & $14{,}3\dB$ & 4,15 & $-1{,}5\dB$ & 12,9\pct{} & 3,8\\
$K=2$ sin compensar & 4 & $23{,}8^\circ$ & $8{,}3\dB$ & 6,27 & $+4{,}5\dB$ & 36\pct{} & 3,0\\
$K=2$ + adelanto & 4 & \multicolumn{6}{c}{\textbf{inestable} (modo elástico)}\\
$K=2$ + atraso & 4 & $75{,}1^\circ$ & $19{,}2\dB$ & 2,52 & $-6{,}7\dB$ & 5,3\pct{} & 10,6\\
\bottomrule
\end{tabular}
\end{center}
El atraso duplica la precisión ante rampa y deja el modo elástico bien por debajo de $0\dB$. A cambio, el lazo es más lento
($\wc$ baja de 4,15 a 2,5) y aparece la cola del dipolo. Moraleja: \textbf{con un modo resonante cerca del cruce, el adelanto lo empeora;
el atraso (o un filtro \emph{notch} centrado en $\omega_f$) lo controla.}

\begin{center}
\begin{tikzpicture}
\begin{axis}[xmode=log,minor x tick num=8,width=0.8\linewidth,height=5.2cm,xmin=0.1,xmax=100,ymin=-50,ymax=40,
  xlabel={$\omega$ [rad/s]},ylabel={$|L|$ [dB]},legend pos=south west,legend columns=2]
\addplot[gris,thick,densely dotted] table[x=w,y=m1]{datos/avion.dat};
\addplot[antes] table[x=w,y=m0]{datos/avion.dat};
\addplot[rojo,thick] table[x=w,y=mlead]{datos/avion.dat};
\addplot[azul,thick] table[x=w,y=mlag]{datos/avion.dat};
\legend{$K=1$,$K=2$,$K=2$ + adelanto,$K=2$ + atraso}
\draw[black] (axis cs:0.1,0)--(axis cs:100,0);
\draw[verde,dashed] (axis cs:0.1,-6)--(axis cs:100,-6);
\node[font=\scriptsize,verde,anchor=south west] at (axis cs:0.11,-6) {límite del pico elástico: $-6\dB$};
\end{axis}
\end{tikzpicture}
\end{center}

% =====================================================================
\appendix
\section{Formulario}
\label{ap:formulario}
{\small
\begin{tcolorbox}[base,colback=white,colframe=azul,title={Bode}]
$|G|_{\mathrm{dB}}=20\log|G|$. Forma de Bode: $K_B=\lim s^NG$. Primer orden: $\pm20\dBdec$, error de $\pm3\dB$ en el quiebre, fase de $\pm45^\circ$ en el quiebre.
Segundo orden: $\pm40\dBdec$, error de $-20\log(2\zeta)$ en $\omega_n$, $M_r=\frac{1}{2\zeta\sqrt{1-\zeta^2}}$ y $\omega_r=\omega_n\sqrt{1-2\zeta^2}$ ($\zeta<0{,}707$).
Retardo: $0\dB$ y fase $-\omega T$ rad. Constantes de error: tipo 0, $20\log K_p$; tipo 1, la asíntota corta $0\dB$ en $K_v$; tipo 2, en $\sqrt{K_a}$.
\end{tcolorbox}
\begin{tcolorbox}[base,colback=white,colframe=azul,title={Nyquist y márgenes}]
$Z=N+P$ (rodeos al $-1$, horario positivo). $\MF=180^\circ+\angle L(j\wc)$ con $|L(j\wc)|=1$.
$\MG=1/|L(j\wf)|$ con $\angle L(j\wf)=-180^\circ$. $K_c=K\cdot10^{\MG_{\mathrm{dB}}/20}$. $\tau_{\max}=\MF[\mathrm{rad}]/\wc$.
Diseño: $\MF$ de $45^\circ$ a $60^\circ$ y $\MG\ge6$--$10\dB$.
\end{tcolorbox}
\begin{tcolorbox}[base,colback=white,colframe=azul,title={Lazo cerrado y correlación (2.º orden)}]
$\zeta\approx\MF/100$. $M_p=e^{-\pi\zeta/\sqrt{1-\zeta^2}}$. $t_s\approx4/(\zeta\omega_n)$. $\wc\lesssim\omega_{BW}\lesssim2\wc$.
$\MF=\arctan\frac{2\zeta}{\sqrt{\sqrt{1+4\zeta^4}-2\zeta^2}}$. $\omega_{BW}=\omega_n\sqrt{1-2\zeta^2+\sqrt{4\zeta^4-4\zeta^2+2}}$.
\end{tcolorbox}
\begin{tcolorbox}[base,colback=white,colframe=verde,title={Adelanto: $G_c=K\frac{Ts+1}{\alpha Ts+1}$\,,\ $\alpha<1$}]
$\phi_m=\MF_{\text{req}}-\MF_0+\delta$ ($\delta$ de $5^\circ$ a $12^\circ$). $\alpha=\frac{1-\sin\phi_m}{1+\sin\phi_m}$.
$\omega_m$: frecuencia donde $|L_0|=-10\log(1/\alpha)$ dB, que es el nuevo $\wc$. $1/T=\omega_m\sqrt\alpha$ y $1/(\alpha T)=\omega_m/\sqrt\alpha$. Verificar.
\end{tcolorbox}
\begin{tcolorbox}[base,colback=white,colframe=verde,title={Atraso: $G_c=K\frac{Ts+1}{\beta Ts+1}$\,,\ $\beta>1$}]
$\wc'$: frecuencia donde $\angle L_0=-180+\MF_{\text{req}}+\delta$ ($\delta\approx6^\circ$). $\beta=10^{|L_0(j\wc')|_{\mathrm{dB}}/20}$.
$1/T=\wc'/10$ y $1/(\beta T)=\wc'/(10\beta)$. Verificar.
\end{tcolorbox}
\begin{tcolorbox}[base,colback=white,colframe=verde,title={Atraso-adelanto}]
Elegir $\wc'$ (por ejemplo, $\wf$). Adelanto centrado en $\wc'$ con $\phi_m=\MF_{\text{req}}-\MF_0(\wc')+6^\circ$.
Atraso con $\beta=10^{A/20}$, donde $A=|L_0(j\wc')|_{\mathrm{dB}}+10\log(1/\alpha)$, y cero en $\wc'/10$.
\end{tcolorbox}
}

\section{Verificación con MATLAB}
\label{ap:matlab}
Las funciones de la \emph{Control System Toolbox} que cubren todo el apunte son:
\begin{center}
\small
\begin{tabular}{@{}ll@{}}
\toprule
\texttt{bode(L)}, \texttt{bode(L1,L2)} & Diagrama de Bode (superpone varios sistemas)\\
\texttt{margin(L)}, \texttt{[Gm,Pm,Wpc,Wgc]=margin(L)} & Bode con márgenes / valores numéricos ($G_m$ en veces)\\
\texttt{allmargin(L)} & Todos los cruces, márgenes de retardo y estabilidad\\
\texttt{nyquist(L)}, \texttt{nichols(L)}; \texttt{ngrid} & Nyquist y Nichols (con contornos M)\\
\texttt{T=feedback(L,1)} & Lazo cerrado con realimentación unitaria\\
\texttt{bandwidth(T)}, \texttt{getPeakGain(T)} & Ancho de banda y $M_r$\\
\texttt{step(T)}, \texttt{stepinfo(T)}, \texttt{lsim(T,u,t)} & Respuestas temporales\\
\texttt{dcgain(minreal(s*L))} & $K_v$\\
\texttt{controlSystemDesigner(G)} & Diseño interactivo arrastrando polos y ceros sobre el Bode\\
\bottomrule
\end{tabular}
\end{center}

Diseño completo del ejemplo~\ref{ej:lead} (adelanto). El archivo \texttt{compensacion.m} reproduce todos los ejemplos:
\begin{tcolorbox}[base,colback=gris!4,colframe=gris,boxrule=0.4pt]
\small
\begin{verbatim}
s  = tf('s');
G  = 1/(s*(s+1)*(0.05*s+1));
K  = 10;  L0 = K*G;                      % 1. ganancia por Kv
[~,MF0,~,wc0] = margin(L0);
phim  = 45 - MF0 + 12;                   % 2. fase a agregar (delta = 12)
alfa  = (1-sind(phim))/(1+sind(phim));   % 3.
magdB = @(x) 20*log10(abs(evalfr(L0,1j*x)));   % 4. |L0| = -10 log(1/alfa)
wm    = fzero(@(x) magdB(x) + 10*log10(1/alfa), [wc0 100]);
T     = 1/(wm*sqrt(alfa));               % 5.
Gc    = K*(T*s+1)/(alfa*T*s+1);
L     = Gc*G;                            % 6. verificar
figure; margin(L); figure; step(feedback(L0,1), feedback(L,1), 8)
\end{verbatim}
\end{tcolorbox}

\end{document}
